%% manuscript_corrections.tex -- master plan of corrections to PAPER_B_MANUSCRIPT.tex.
%% GENERATED by notes/tools/make_corrections_tex.py from notes/tools/corrections_entries.py
%% and notes/tools/plan_entries_E.py (BEFORE text cut from the committed manuscript by
%% script). Edit the entries files, then run:
%%   python3 notes/tools/make_corrections_tex.py notes/manuscript_corrections.tex
\documentclass[10pt]{article}
\usepackage[a4paper,landscape,margin=1.4cm]{geometry}
\usepackage{amsmath,amssymb,dsfont}
\usepackage{amsthm}
\newtheorem{proposition}{Proposition}
\newtheorem{lemma}{Lemma}
\newtheorem{theorem}{Theorem}
\theoremstyle{definition}
\newtheorem{definition}{Definition}
\usepackage{longtable,array,booktabs}
\usepackage[authoryear,round]{natbib}
\usepackage[T1]{fontenc}
\newcommand{\PB}{P^{\mathrm{B}}}
\newcommand{\PJ}{P^{\mathrm{J}}}
\newcommand{\Pbar}{\overline{P}}
\newcommand{\assoc}{\operatorname{assoc}}
\newcommand{\sgn}{\operatorname{sgn}}
\newcommand{\E}{\mathbb{E}}
\newcommand{\bigO}{\mathcal{O}}
\newcommand{\vv}{v}
\newcommand{\anchor}[1]{}
\renewcommand{\ref}[1]{\textsc{\detokenize{#1}}}
\renewcommand{\footnote}[1]{ \textit{[Footnote: #1]}}
\newlength{\fullw}\setlength{\fullw}{25.9cm}
\newcommand{\entryhead}[2]{\multicolumn{3}{|l|}{\rule{0pt}{2.6ex}\large\textbf{#1}\quad #2} \\ \hline}
\setlength{\parindent}{0pt}
\title{Paper B: corrections to the manuscript (master plan)}
\author{}
\date{Against the manuscript at db3a2f41, 2 October 2026}
\begin{document}
\maketitle
\vspace{-1.5em}
\noindent Each row sets the current manuscript text (left) against the proposed replacement (right),
both typeset as they would appear. BEFORE text is cut from the committed manuscript by script, so it
matches the manuscript exactly. Under each entry, \textbf{Why} gives the error and \textbf{Evidence} the
Lean or sympy record or page that settles it (audit items refer to \texttt{notes/citation\_audit.md}).
The identification paragraph appended in 3.4 reconciles Section 3 with
\texttt{notes/positioning\_economics.tex}. Cross-references print as their label (e.g.\ \textsc{prop:IMM});
footnotes print inline. \textbf{Nothing here is applied.} The author approves entries by number;
approved entries are applied to the manuscript, compiled and committed. Every AFTER text passes the
mechanical checks of \texttt{notes/writing\_discipline.md} (no colons, no dashes doing a sentence's
work, ``sequence'' for reading order, the fixed adoption terminology, length within 80--120\% of the
draft except where the entry says why); the generator refuses to build otherwise.

\section*{Entries in manuscript order}
Entries are numbered by the manuscript section they change and listed in the order in which their
text would appear. \textbf{0} is the abstract, \textbf{1} to \textbf{7} are the numbered sections,
\textbf{A} is the appendix and \textbf{B} is the back matter with the bibliography. Each entry also
shows its former number, which earlier notes and commits use, and the table at the end maps old
numbers to new. Order of application, where it differs from manuscript order. B.2 (bibliography) goes in before every entry that cites a new reference. 5.2 (Proposition ORD) and 6.3 (Propositions ADJ, LAD and FAC) go in before 1.3, 1.4, 2.3 and 1.9, which cite them. 6.2 goes in before 6.3, which inserts after it. B.1 goes in last.

\paragraph*{Applied to the manuscript} at db3a2f41: 2.1, 2.2, 2.3, 2.4, 2.5, 3.1, 3.4, 4.1, 5.1, 5.2, 6.1, 6.2, 6.3.
These entries appear in the table of numbers at the end and no longer in the table of changes.

\paragraph*{Pending} (18 entries, in manuscript order): 0.1, 1.1, 1.2, 1.3, 1.4, 1.5, 1.6, 1.7, 1.8, 1.9, 2.6, 3.2, 3.3, 3.5, 4.2, A.1, B.1, B.2.

{\small
\begin{longtable}{|p{0.9cm}|p{12.2cm}|p{12.2cm}|}
\hline
\textbf{Part} & \textbf{BEFORE (current manuscript, or the anchor for an insert)} & \textbf{AFTER (proposed)} \\ \hline\hline
\endhead
\multicolumn{3}{|l|}{\rule{0pt}{3.2ex}\Large\textbf{0\quad Abstract}} \\ \hline\hline
\entryhead{0.1}{Abstract, sentences 1, 2-4 and 6 (rebased on the author's draft c2ae782f)\quad{\footnotesize\textit{(formerly C.1)}}}
1/3 & The effect of the order in which evidence arrives on an individual judgment has sparked much interest in both experimental and theoretical studies. & The effect of the sequence in which evidence arrives on an individual judgment has sparked much interest in both experimental and theoretical studies. \\ \hline
2/3 & Considering a case where evaluators make decisions using two correlated soft cues, the current papers asks which of their beliefs and decision statistics register a resulting effect of arrival sequence with respective to a sequence-free benchmark. With a view of implications for empirical measurement, the paper shows that an arbitrary statistic is an order of magnitude closer to the benchmark when it is insensitive to small shifts in the marginal probabilities. The paper shows how the collective marginal probabilities and the share of evaluators whose decision the sequence changes depart from the sequence-free benchmark at first order, whereas the believed association departs only at second order. & Considering a case where evaluators make decisions using two correlated soft cues, the current paper asks which statistics of their beliefs and decisions register a resulting effect of arrival sequence with respect to a sequence-free benchmark. With a view to implications for empirical measurement, the paper shows that a smooth statistic of the belief is an order of magnitude closer to the benchmark than the marginal probabilities when it is insensitive to small shifts in those marginals. The average marginal probabilities and the share of evaluators whose decision the sequence changes depart from the sequence-free benchmark at first order, whereas the believed association departs only at second order. \\ \hline
3/3 & Second, whether the effect of order is detectable or not depends on the kind of question being asked rather than on the number of evaluators considered. & Second, at a given precision of measurement, whether the sequence effect is detectable depends on the kind of question being asked. \\ \hline
\multicolumn{3}{|p{\fullw}|}{\textbf{Why.} The author applied 0.1 in their own wording at c2ae782f; this entry keeps only what still needs correcting. Sentence 1 and the last sentence use ``order'' for reading sequence, which the paper reserves for order in $c$ (writing discipline 6). Sentence 2: ``the current papers asks'', ``with respective to'', and ``which of their beliefs and decision statistics'' (a belief is not a statistic). Sentence 3: ``With a view of'' is ``With a view to''; ``an arbitrary statistic'' is false as stated, since only statistics insensitive to the marginals are closer (Proposition PRO), so ``smooth ... of the belief'', the comparator ``than the marginal probabilities'' and ``those marginals'' are restored from the approved 0.1. Sentence 4: two consecutive sentences open ``The paper shows'', and ``shows how'' states a result as a method; ``collective'' becomes ``average'' (Proposition DRF reads the mean belief). Leaving the loss out of sentence 4 is accurate, since the loss is also second order (Theorem LOS). Last sentence: ``rather than on the number of evaluators'' is the identification overclaim (todo, 2026-10-01). Only the odds ratio is blind at every sample size (Lemma SEP); the association is second order and visible once precision is finer than $c^2$ (MS precision passage, $c^2\ll\varepsilon\ll c$), and precision improves with the number of evaluators. Fixing the precision makes the claim true.} \\
\multicolumn{3}{|p{\fullw}|}{\textbf{Evidence.} PRO (propPRO\_protection, propPRO\_uniqueness); DEC and LOS (second order); DRF and SHR (first order); LemmaSEP.lean (the odds ratio). All in the Lean and sympy suites.} \\ \hline
\multicolumn{3}{|p{\fullw}|}{\textbf{Grounds.}\par \textbf{Theorem.} \textit{PropPRO.lean, propPRO\_protection} \newline For all real $x,y$ and all $(\alpha,\beta,\lambda,q_0,r_0)$,
$\langle xJ + y\,\nabla\mathrm{assoc}(q\otimes r),\; M_\lambda\rangle = 0$, where
$M_\lambda = \lambda\kappa R_1 + (1-\lambda)\kappa' R_2$, $R_1 = q\otimes(1,-1)$,
$R_2 = (1,-1)\otimes r$, $\kappa = (\alpha-q_0)r_0(1-r_0)/Z$,
$\kappa' = (\beta-r_0)q_0(1-q_0)/Z$ and $Z = \alpha\beta(1-\alpha)(1-\beta)$. \newline \textit{Proposition PRO (i): a statistic whose gradient lies in the plane
$\mathrm{span}\{J,\nabla\mathrm{assoc}\}$ has zero first-order coefficient,
identically in the prior and in the mix.}} \\
\multicolumn{3}{|p{\fullw}|}{\textbf{Theorem.} \textit{PropPRO.lean, propPRO\_uniqueness and annihilator\_eq\_span} \newline Let $\lambda\in(0,1)$, $q_0,r_0\in(0,1)$, $\beta\neq r_0$ and
$x = \langle\nabla F,R_1\rangle$, $y = \langle\nabla F,R_2\rangle$. If
$\lambda(\alpha-q_0)r_0(1-r_0)\,x + (1-\lambda)(\beta-r_0)q_0(1-q_0)\,y = 0$ at two
priors $\alpha_1\neq\alpha_2$, then $x = y = 0$; and any $V$ with
$\langle V,R_1\rangle = \langle V,R_2\rangle = 0$ equals
$(V_{01} + r_0(V_{00}-V_{01}))J + \frac{V_{00}-V_{01}}{1-q_0}\nabla\mathrm{assoc}$. \newline \textit{Proposition PRO (ii): protection on an open set of priors forces
$\nabla F$ into the plane, so the classification is an ``if and only if''.}} \\
\multicolumn{3}{|p{\fullw}|}{\textbf{Theorem.} \textit{PropPRO.lean, margA\_not\_in\_span; Aggregate.lean, propDRF} \newline For $0<r_0<1$ the $A$-marginal read-out $(0,0,1,1)$ is not of the
form $xJ + y\,\nabla\mathrm{assoc}(q_0,r_0)$. For $\alpha,\beta\notin\{0,1\}$ and
$Z\neq0$,
$\frac{d}{dc}\big[\lambda(P^J_{AB}-P^B)(A{=}1) + (1-\lambda)(P^J_{BA}-P^B)(A{=}1)\big]_{c=0}
= (1-\lambda)\,q_0(1-q_0)(r_0-\beta)/Z$. \newline \textit{Proposition DRF: a marginal departs at first order for every
$\lambda<1$, so ``an arbitrary statistic'' is not an order closer; only the
protected ones are.}} \\
\multicolumn{3}{|p{\fullw}|}{\textbf{Theorem.} \textit{PropPRO.lean, lemmaASC\_first\_order\_vanishes (and \_BA)} \newline For all $q_0,r_0,t,\kappa$:
$\mathrm{assoc}\big(q\otimes r + t\kappa R_1\big) = \mathrm{assoc}(q\otimes r)
+ t^2\,\mathrm{assoc}(\kappa R_1)$, and likewise along $R_2$: the first-order
term $t\,\langle\nabla\mathrm{assoc},\kappa R_\sigma\rangle$ is identically zero. \newline \textit{Proposition DEC: the believed association is in the protected
class, second order along either route.}} \\
\multicolumn{3}{|p{\fullw}|}{\textbf{Theorem.} \textit{Decision.lean, volume\_flipSet and lintegral\_stake\_le\_volume} \newline With $u$ the benchmark surplus, $\delta$ the first-order score
coefficient and $\mathrm{Flips}(u,c,\delta) := \neg\big(0\le u \iff 0\le u + c\delta\big)$:
$\mathrm{Leb}\{u : \mathrm{Flips}(u,c,\delta)\} = |c\delta|$ and
$\int_{\{\mathrm{Flips}\}} |u|\,du \le |c\delta|^2$, for all real $c,\delta$. \newline \textit{Proposition SHR and Theorem LOS: the affected share is first order,
the surplus-weighted loss second order.}} \\ \hline
\multicolumn{3}{|l|}{} \\ \hline
\multicolumn{3}{|l|}{\rule{0pt}{3.2ex}\Large\textbf{1\quad Introduction}} \\ \hline\hline
\entryhead{1.1}{Introduction, paragraph 1 (rebased on the author's draft c2ae782f)\quad{\footnotesize\textit{(formerly C.2)}}}
1/5 & That arrival sequence of cues does effect judgment is rarely under doubt -- since if individual judgments did not depend on the arrival sequence, then each cue must move belief by a factor fixed before it is received and lead to the unlikely situation where experience does not change how later cues are read at all. & That the arrival sequence of cues affects judgment is rarely in doubt, since if individual judgments did not depend on the arrival sequence, then each cue must move belief by a factor fixed before it is received, leading to the unlikely situation where experience does not change how later cues are read at all. \\ \hline
2/5 & The sequence-dependence of Jeffrey conditioning for soft evidence in particular has attracted much debate \citep{DiaconisZabell1982,Hawthorne2004}. However, but the question of how that dependence registers in the statistics an observer reads has received little attention. & The sequence dependence of Jeffrey conditioning for soft evidence in particular has attracted much debate \citep{DiaconisZabell1982,Hawthorne2004}, but the question of how that dependence registers in the statistics an observer reads has received little attention. \\ \hline
3/5 & It finds that the share of population affected by evidence sequence deviates from a sequence-invariant benchmark -- but the believed association between attributes and the total loss due to sequence effects do not. & It finds that the share of the population whose decision the sequence changes departs from a sequence-invariant benchmark at first order, while the believed association between attributes and the surplus-weighted loss depart from it only at second order. \\ \hline
4/5 & Given that a difference from the benchmark below the observer's precision is invisible in the statistic read, the paper discusses the implications for identification of sequence effects. & \textit{(sentence deleted)} \\ \hline
5/5 & Despite its simplicity the two-cue setting demonstrates the central mechanism of how an attribute-local updating leaves the interactions among attributes untouched remains valid for any number of cues and attributes. & Despite its simplicity, the two-cue setting demonstrates the central mechanism, that attribute-local updating leaves the interactions among attributes untouched, and the mechanism remains valid for any number of cues and attributes (Lemma~\ref{lem:SEP}). \\ \hline
\multicolumn{3}{|p{\fullw}|}{\textbf{Why.} The author applied 1.1 in their own wording at c2ae782f; this entry keeps what still needs correcting. Modus tollens sentence: ``does effect'' is ``affects'', ``under doubt'' is ``in doubt'', the double hyphen is a dash doing a sentence's work, and ``and lead to'' breaks the parallel with ``move''. Caution kept from the earlier entry: ``must'' is stronger than the literature. Fixed factors imply commutation (Wagner 2003 Theorem 3.2, Hawthorne factorUpdate\_comm, Cripps order\_invariance); the converse is exact in this paper's two-cue model, where the sequence effect vanishes iff q0 = alpha and r0 = beta, but is not a theorem for arbitrary updating rules. Debate sentence: ``However, but'' and the hyphen in the noun ``sequence dependence''. ZhaoOsherson2010 is no longer cited anywhere in the manuscript, which is the author's choice; its bibliography entry is now unused. Sentence ``It finds that'': the loss is not free of the sequence effect, it carries it at second order (LOS), and ``-- but'' is a dash doing a sentence's work. Sentence ``Given that'': cut (rule 4 of the writing discipline); the precision point is made in the 1.3 paragraph, and ``identification of sequence effects'' is the identification overclaim 0.1 removes from the abstract. Last sentence: ``demonstrates the central mechanism of how ... remains valid'' has two verbs for one subject; Lemma SEP is the general-N result it refers to.} \\
\multicolumn{3}{|p{\fullw}|}{\textbf{Evidence.} verify\_jeffrey.md J1; verify\_kinematics.md; verify\_hawthorne\_weisberg.md; Theorem LOS, Proposition SHR; LemmaSEP.lean (general N). Modus tollens: Cripps.lean order\_invariance, Hawthorne.lean factorUpdate\_comm, Wagner2003.lean thm32; sympy/verify\_soft\_vs\_hard.py (exact condition for the sequence effect).} \\ \hline
\multicolumn{3}{|p{\fullw}|}{\textbf{Grounds.}\par \textbf{Quote.} \textit{Zhao and Osherson 2010, p. 290 (opening part, before Experiment 1)} \newline ``The normative status of Jeffrey's rule has nonetheless been
questioned because successive uses produce distinct distributions depending on
the order in which events are considered (D\"{o}ring, 1999). In our view, such
doubts disappear on closer inspection of the evidential weight of probability
judgements (Osherson, 2002; Wagner, 2002).'' \newline \textit{They report the doubt as settled, not the mechanism as ``unclear''.}} \\
\multicolumn{3}{|p{\fullw}|}{\textbf{Quote.} \textit{Jeffrey 1983, pp. 182-183 (Section 11.11)} \newline ``It is straightforward to verify that the present kinematical
scheme is not generally commutative \ldots the resulting probability measure
may be one thing or another depending on the order of the two applications.
That is as it should be, for just after the $A$-application, one's degrees of
belief in the $A_i$ \emph{are} $a_i$, no matter what they were before \ldots
The notion that the kinematical scheme \emph{ought} to be generally commutative
(Domotor, \emph{Philosophy of Science} 47 [1980]: 395) stems from a conflation
of two attitudes toward The Given'' \newline \textit{Jeffrey holds non-commutativity to be correct and the demand for
commutativity to rest on a conflation.}} \\
\multicolumn{3}{|p{\fullw}|}{\textbf{Quote.} \textit{Diaconis and Zabell 1982, p. 827, Section 4.2, Remark 2} \newline ``Thus noncommutativity is not a real problem for successive
Jeffrey updating.'' \newline \textit{D-Z are on the side of the sequence effect, not of the doubt.}} \\
\multicolumn{3}{|p{\fullw}|}{\textbf{Quote.} \textit{Hawthorne 2004, p. 99 (Section 5)} \newline ``The second concern is that, as a result of this amnesia, the
order in which states or experiences are acquired will almost always have a
very significant influence on an agent's belief strengths. This order effect
is very troubling.'' \newline \textit{The one cited author who treats the order effect as a defect; the
three sources disagree, which is a debate.}} \\
\multicolumn{3}{|p{\fullw}|}{\textbf{Theorem.} \textit{Decision.lean, lintegral\_stake\_le\_volume with volume\_flipSet} \newline For all real $c,\delta$:
$\int_{\{u:\mathrm{Flips}(u,c,\delta)\}} |u|\,du \le |c\delta|\cdot|c\delta|$,
whereas $\mathrm{Leb}\{u:\mathrm{Flips}(u,c,\delta)\} = |c\delta|$. \newline \textit{Theorem LOS: the loss carries the sequence effect at second order;
it is not free of it, and the share (Proposition SHR) is first order.}} \\
\multicolumn{3}{|p{\fullw}|}{\textbf{Theorem.} \textit{Hawthorne.lean, factorUpdate\_comm; Cripps.lean, order\_invariance} \newline With $\mathrm{factorUpdate}(p,u,w)(x)=w(u(x))\,p(x)/\sum_y w(u(y))\,p(y)$, a
belief multiplied by fixed factors $w$ on one partition and $z$ on another, both
normalisers nonzero:
$\mathrm{factorUpdate}(\mathrm{factorUpdate}(p,u,w),v,z)
=\mathrm{factorUpdate}(\mathrm{factorUpdate}(p,v,z),u,w)$. Likewise, under Cripps's
Symmetry and Divisibility, two signals with fixed likelihoods give the same posterior
in either sequence. \newline \textit{(sentence 1) The first half of the modus tollens. Read as Bayes factors,
fixed before either cue arrives, the cues cannot leave a trace of the sequence.}} \\
\multicolumn{3}{|p{\fullw}|}{\textbf{Computation.} \textit{sympy/verify\_soft\_vs\_hard.py, symbolic checks (15/15)} \newline Symbolic in $\alpha,\beta,c,q_0,r_0$: a credential read first moves the
$B$-marginal the letter meets to $\beta+c(q_0-\alpha)/(\alpha(1-\alpha))$, so the
factor the letter implies changes unless $c=0$ or $q_0=\alpha$, and symmetrically for
the credential. Every cell of $\PJ_{AB}-\PJ_{BA}$ carries the factor $c$, and with
$c\neq0$ the difference vanishes if and only if $q_0=\alpha$ and $r_0=\beta$. \newline \textit{(sentence 1) Under the delivered-credence reading the sequence leaves a
trace exactly when an earlier cue changes the factor a later one implies, which is
the lived-experience point in the model's terms. The sentence itself claims only the
first half; this row shows that the delivered-credence reading, one of the readings it
leaves open, is one in which what came earlier changes how later cues are read.}} \\ \hline
\multicolumn{3}{|l|}{} \\ \hline
\entryhead{1.2}{Introduction, paragraph 2 and its footnote (rebased on the author's draft c2ae782f)\quad{\footnotesize\textit{(formerly C.3)}}}
 & The discussions on how Jeffrey\citep{Jeffrey1983} conditioning depends on arrival sequence have focused more often on whether the commutation criterion is exactly satisfied or not but not enough on how measurable its failure is in observed data. The belief-adjustment literature often models an impression as a single score \citep{Hawthorne2004} or considers traits one attribute at a time \citep{Asch1946}. The characterisation of when two sequences lead to the same belief does consider a joint belief over attributes\citet{DiaconisZabell1982}, but the order dependence is more more treated as a defect to repair \footnote{\citet{Hawthorne2004} repairs this by letting a cue multiply the old belief instead of replacing it, and under his order-free variants the sequence effect disappears.} rather than as something to measure. While changing the input from a credence that replaces the prior to a factor that multiplies it -- does settle whether beliefs depend on sequence, it hardly addresses the empirical issue that an observer cannot tell whether impressions replace prior belief or adjust it. & The discussions on how Jeffrey conditioning \citep{Jeffrey1983} depends on arrival sequence have focused more often on whether the commutation criterion is exactly satisfied or not but not enough on how measurable its failure is in observed data. The belief-adjustment literature often models an impression as a single score \citep{HogarthEinhorn1992}, and the impression-formation literature considers traits one attribute at a time \citep{Asch1946}. The characterisation of when two sequences lead to the same belief does consider a joint belief over attributes \citep{DiaconisZabell1982}, but elsewhere sequence dependence is more often treated as a defect to repair\footnote{\citet{Hawthorne2004} offers factor-based alternatives in which a cue supplies a normed-likelihood or likelihood-ratio factor instead of a new probability. Under these, updates on distinct attributes commute, and in his Basis-Commuting Version all updates do.} rather than as something to measure. While changing the input from a credence that replaces the prior to a factor that multiplies it does settle whether beliefs depend on sequence, it hardly addresses the empirical issue that an observer cannot tell whether impressions replace prior belief or adjust it. \\ \hline
\multicolumn{3}{|p{\fullw}|}{\textbf{Why.} The author applied 1.2 in their own wording at c2ae782f; this entry keeps what still needs correcting. (i) ``models an impression as a single score \textbackslash{}citep\{Hawthorne2004\}'': the single-score model is Hogarth-Einhorn's belief-adjustment model, and Hawthorne is not part of that literature. (ii) Asch is impression formation, not belief adjustment, and reports eighteen traits one by one (audit A28). (iii) ``\textbackslash{}citet\{DiaconisZabell1982\}'' is glued to ``attributes'' and should be \textbackslash{}citep. (iv) ``the order dependence is more more treated as a defect'': doubled ``more'', ``order'' for reading sequence, and the clause still groups Diaconis-Zabell with the defect view, though they hold that non-commutativity ``is not a real problem'' (M24); ``elsewhere'' detaches it from them. (v) The footnote still says Hawthorne ``repairs'' the defect and that the effect disappears; he offers alternatives, and freedom from sequence holds across distinct bases, fully only in his Basis-Commuting Version (M19). (vi) ``Jeffrey\textbackslash{}citep'' lacks a space and the citation sits inside ``Jeffrey conditioning''. (vii) The double hyphen before ``does settle'' is a dash doing a comma's work.} \\
\multicolumn{3}{|p{\fullw}|}{\textbf{Evidence.} Asch1946 record (Asch.lean, Table 7); HogarthEinhorn.lean; DiaconisZabell.lean and README; Hawthorne.lean (\texttt{extUpdate\_basisCommuting}, \texttt{factorUpdate\_comm}), Sections 6-8.} \\ \hline
\multicolumn{3}{|p{\fullw}|}{\textbf{Grounds.}\par \textbf{Theorem.} \textit{Asch.lean, seriesA\_length, seriesB\_eq\_reverse, checkListI\_length, allTraits\_length, table7\_rows\_are\_checklist\_pairs, d\_ne\_zero, restrained\_64\_9} \newline $|\mathrm{Series\,A}| = 6$ and $\mathrm{Series\,B} = \mathrm{reverse}(\mathrm{Series\,A})$;
$|\mathrm{Check\,List\,I}| = 18$ and Table 7 has $18$ rows, row $i$ naming a member
of pair $i$; for every one of the $18$ traits
$\mathrm{ie}(t) - \mathrm{ei}(t) \neq 0$; e.g.\ restrained $64$ vs $9$,
good-looking $74$ vs $35$, humorous $52$ vs $21$, good-natured $18$ vs $0$. \newline \textit{Experiment VI reports eighteen per-trait percentages under two
orders, i.e.\ eighteen marginals, not a single score.}} \\
\multicolumn{3}{|p{\fullw}|}{\textbf{Quote.} \textit{Asch 1946, pp. 270-271, Experiment VI and Table 7} \newline p.~270: ``The two series are identical with regard to their
members, differing only in the order of succession of the latter.''
p.~271: ``The check-list data appearing in Table 7 furnish quantitative
support for the conclusions drawn from the written sketches.'' \newline \textit{The check list is Asch's own quantitative instrument for the
order comparison.}} \\
\multicolumn{3}{|p{\fullw}|}{\textbf{Quote.} \textit{Diaconis and Zabell 1982, p. 827, Section 4.2, Remarks 1-2} \newline ``Remark 1. There is no reason to require
$P_{\mathcal{E}\mathcal{F}} = P_{\mathcal{F}\mathcal{E}}$ for successive
updating to be useful and valid. \ldots Remark 2. \ldots Thus noncommutativity
is not a real problem for successive Jeffrey updating.'' \newline \textit{D-Z do not hold the defect view the manuscript groups them with.}} \\
\multicolumn{3}{|p{\fullw}|}{\textbf{Quote.} \textit{Hawthorne 2004, pp. 96 and 108 (Sections 4 and 8)} \newline p.~96: ``Indeed, it may turn out that there is no one true theory of
uncertain updating -- that each theory has its uses, its domain of
applicability.'' p.~108: ``Let us call this the Basis-Commuting Version of the
Likelihood-Ratio Update Model. Sequential updating is completely
order-independent on this model.'' \newline \textit{He offers alternatives rather than a repair, and full order-freedom
is claimed only for the Basis-Commuting Version.}} \\
\multicolumn{3}{|p{\fullw}|}{\textbf{Theorem.} \textit{Hawthorne.lean, factorUpdate\_comm and extUpdate\_basisCommuting} \newline For bases $u:\Omega\to\iota$, $v:\Omega\to\kappa$ and factor vectors
$w,z$ with $\sum_y w(u(y))p(y)\neq0$ and $\sum_y z(v(y))p(y)\neq0$:
$\mathrm{fU}(\mathrm{fU}(p,u,w),v,z) = \mathrm{fU}(\mathrm{fU}(p,v,z),u,w)$. If the
factor $\Lambda_k(\gamma)$ of every basis-homogeneous subsequence is invariant
under permutation of $\gamma$, then $\mathrm{extUpdate}(p,\delta) =
\mathrm{extUpdate}(p,\delta')$ for every permutation $\delta'$ of $\delta$. \newline \textit{Factor updates commute across distinct bases; commutation within a
basis, hence full order-freedom, needs the extra Basis-Commuting hypothesis.}} \\ \hline
\multicolumn{3}{|l|}{} \\ \hline
\entryhead{1.3}{Introduction, the two-channel paragraph (rebased on the author's draft c2ae782f)\quad{\footnotesize\textit{(formerly C.4)}}}
 & Whether impressions replace prior belief or adjust is a question which \citet{Hawthorne2004}, writing as a logician rather than a psychologist, leave open. To set the scope of its conclusions, the paper considers that an arrival-sequence dependence can enter through two channels, an \textbf{association} channel, in which one cue changes what the other implies and which exists only when the attributes are believed related, and a \textbf{position} channel, in which the observer weights the later cue less whatever the attributes are \citep{HogarthEinhorn1992,Asch1946}. In the position channel, the later cue is adopted only in part and every marginal registers the sequence even when the attributes are believed unrelated. On the other hand, when each impression is adopted in full, as in the successive updating of \citet{DiaconisZabell1982}, the position channel is disabled with only the association channel remaining. In the two-attribute setting with full-adoption, the paper thus considers three coordinates of a belief -- two marginal probabilities and the cross-attribute association -- and explains how a difference from the sequence-free benchmark below the observer's precision is invisible in the statistic read by the observer. The question is therefore not simply whether the two sequences agree but by how much they disagree in an arbitrary statistic. The paper finds that the marginal probabilities and decisions from them carry the difference at first order in the prior covariance, while the believed association and statistics carry it only at second order. This carries clear implications for what an audit can or cannot measure about sequence dependence. & Whether impressions replace prior belief or adjust it is a question which \citet[pp.~115--116]{Hawthorne2004}, writing as a logician rather than a psychologist, leaves open. To set the scope of its conclusions, the paper considers that arrival-sequence dependence can enter through two channels, an \textbf{association} channel, in which one cue changes what the other implies and which exists only when the attributes are believed related, and a \textbf{position} channel, in which the weight a cue receives depends on where in the sequence it arrives, whatever the attributes are \citep{HogarthEinhorn1992}. In the position channel the later cue is adopted only in part and every marginal registers the sequence even when the attributes are believed unrelated. On the other hand, when each impression is adopted in full, as in the successive updating of \citet{DiaconisZabell1982}, the position channel is disabled with only the association channel remaining. In the two-attribute setting with full adoption, the paper thus considers three coordinates of a belief, the two marginal probabilities and the cross-attribute association, and explains how a difference from the sequence-free benchmark below the observer's precision is invisible in the statistic read by the observer. The question is therefore not simply whether the two sequences agree but by how much they disagree in an arbitrary statistic. The paper finds that the marginal probabilities and the share of decisions they change carry the difference at first order in the prior covariance, while the believed association and the statistics that move with it carry it only at second order. This carries clear implications for what an audit can or cannot measure about sequence dependence. \\ \hline
\multicolumn{3}{|p{\fullw}|}{\textbf{Why.} The author applied 1.3 in their own wording at c2ae782f; this entry keeps what still needs correcting and adds one sentence. Corrections: ``adjust'' to ``adjust it''; ``leave open'' to ``leaves open'' (Hawthorne is one author; he gives a tentative view and says his interest is normative, pp. 115-116); ``an arrival-sequence dependence'' loses its article; the position channel was defined as ``the observer weights the later cue less'' and cited to Hogarth-Einhorn and Asch, but Hogarth-Einhorn's step-by-step adjustment predicts recency (Appendix B) and Asch denies that ``sheer temporal position'' matters (p. 272), so the definition is made neutral and Asch is dropped (audit P10-P12); the two dashes become commas; ``decisions from them'' becomes ``the share of decisions they change'', since the loss is a decision statistic and is second order (LOS); ``and statistics'' becomes ``and the statistics that move with it''. Scope defence (author, 2026-10-01): the position channel is a property of the evaluator, so the paper sets it aside and studies the measurement problem that remains under the association channel. The author placed the one-sentence version of this in the panel paragraph (1.4) and the full version with its test in Section 6 (6.3), and asked for less in the introduction; this paragraph therefore keeps the author's own two sentences on the two settings and adds nothing.} \\
\multicolumn{3}{|p{\fullw}|}{\textbf{Evidence.} HogarthEinhorn.lean (\texttt{appB\_recency}, \texttt{eq8\_estimation\_first\_dominates}); Hawthorne 2004 pp. 115-116 (grounds\_E\_literature); Ladder.lean (\texttt{ladder\_gap}, \texttt{ladder\_assoc\_coeff}, \texttt{ladder\_oddsShadow\_seqEffect}, \texttt{oddsRatio\_rescale}); sympy/verify\_ladder.py (37/37).} \\ \hline
\multicolumn{3}{|p{\fullw}|}{\textbf{Grounds.}\par \textbf{Theorem.} \textit{HogarthEinhorn.lean, appB\_recency, appB\_recency\_pos, oneSided\_eq\_eq8, oneSided\_primacy\_iff} \newline With $\mathrm{est}(w,S,s) = S + w(s-S)$ (HE Eq.~3/4): for all reals,
$\mathrm{est}(w_b,\mathrm{est}(w_a,S_0,a),b) - \mathrm{est}(w_a,\mathrm{est}(w_b,S_0,b),a)
= w_a w_b (b-a)$, which is $>0$ when $w_a,w_b>0$ and $a<b$ (recency). The
one-sided rule $\mathrm{est}(w,\mathrm{est}(1,S_0,a),b)$ equals HE's Eq.~8 with
$k=2$, and for $a<b$ it gives $S_{ab}-S_{ba}<0$ iff $w<1/2$. \newline \textit{Damping every cue gives recency; primacy comes from adopting the
first cue in full (Eq.~8), and then only for $w<1/2$.}} \\
\multicolumn{3}{|p{\fullw}|}{\textbf{Quote.} \textit{Hogarth and Einhorn 1992, transcription T-p. 35 (Appendix B, Eq. B.5) and T-p. 12 (Eq. 8)} \newline T-p.~35: ``$D = w_a w_b[s(x_b) - s(x_a)]$. Because $s(x_b) > s(x_a)$,
$D > 0$ and recency always obtains.'' T-p.~12: ``$S_k = s(x_1) + w_k[s(x_2,\ldots,x_k) - R]$
\ldots the effective weight accorded to $s(x_1)$ must be greater than the weight
attached to any of the other pieces of evidence.'' \newline \textit{Their step-by-step partial adjustment predicts recency; their
primacy is the End-of-Sequence anchoring on the first item.}} \\
\multicolumn{3}{|p{\fullw}|}{\textbf{Quote.} \textit{Asch 1946, p. 272} \newline ``It is not the sheer temporal position of the item which is
important as much as the functional relation of its content to the content of
the items following it.'' \newline \textit{Asch denies a position channel, so he cannot be cited for ``weights
the later cue less''.}} \\
\multicolumn{3}{|p{\fullw}|}{\textbf{Theorem.} \textit{BenjaminBodohCreedRabin.lean, margOddsA\_orders\_ne and example\_marginal\_gap} \newline On the $2\times2$ joint with prior $u\otimes v$ (independent, $c=0$),
$\alpha\neq1$, positive cue likelihoods $a$ on $A$ and $b$ on $B$, the $A$
cue informative ($a_0\neq a_1$): the odds $P(A{=}0)/P(A{=}1)$ after base-rate-neglect updating
differ between the orders $AB$ and $BA$. Instance: uniform prior,
$\alpha=1/2$, $a=b=(49/50,1/50)$ gives $P(A{=}0)=7/8$ reading $A$ first and
$49/50$ reading $A$ last. \newline \textit{A position channel operates with likelihood inputs adopted in full,
so ``which channels operate is fixed by adoption'' holds for this model only.}} \\
\multicolumn{3}{|p{\fullw}|}{\textbf{Quote.} \textit{Diaconis and Zabell 1982, p. 825 (Section 3.1) and p. 827 (Section 4.2)} \newline p.~825: ``To use Jeffrey's rule at the second stage we must, of
course, accept the J-condition \ldots Clearly, the order of updating matters,
since the second opinion dominates.'' p.~827: ``2. When is successive updating
reasonable?'' \newline \textit{Successive updating with each impression adopted in full is D-Z's
setting, examined, not a premise they lay down.}} \\ \hline
\multicolumn{3}{|l|}{} \\ \hline
\entryhead{1.4}{Introduction, hiring-panel paragraph, last three sentences (rebased on the author's draft d48dd872)\quad{\footnotesize\textit{(formerly C.5)}}}
 & As the paper shows, this means that the believed association differs between the two sequences only at second order. On the other hand, a letter adopted only in part moves that belief from wherever the credential left it (and the believed association differs at first order - see Proposition LAD). Focusing on situations where unrelated traits do not exhibit sequence depends through the positional channel, the focuses on the measurement issues under association channel. & The believed association then differs between the two sequences only at second order (Proposition~\ref{prop:ORD}). A letter adopted only in part moves that belief from wherever the credential left it, and the believed association then differs at first order (Proposition~\ref{prop:LAD}). Since a partly adopted letter makes even unrelated traits sequence dependent through the position channel, the paper sets that channel aside and studies the measurement problem that remains under the association channel. \\ \hline
\multicolumn{3}{|p{\fullw}|}{\textbf{Why.} The author applied the panel paragraph's new ending at d48dd872 in the pre-correction wording; this entry keeps the corrections that remain. ``As the paper shows, this means that'' has no noun for ``this'' (rule 3); the sentence now follows on from the previous one, which ends ``the believed link between the traits''. The hyphen in ``at first order - see Proposition LAD'' is a dash doing a sentence's work (rule 9), and ``Proposition LAD'' needs the reference. ``On the other hand'' goes, the contrast being carried by ``adopted in full'' against ``adopted only in part''. The last sentence is repaired: ``sequence depends'' to ``sequence dependent'', ``the focuses'' to ``the paper sets that channel aside and studies'', ``positional channel'' to the agreed ``position channel'', ``under association channel'' to ``under the association channel''. It is the one place in the introduction that states the scope defence (author, 2026-10-01): the position channel is the evaluator's, so the paper sets it aside and studies the measurement problem that remains. The odds-ratio sentence is in a comment at the author's choice (less in the introduction); its claim, exact invariance at every weight, is Lemma SEP and stays in Section 6. Economics behind the three sentences: under full adoption the letter fixes its own marginal whatever the credential implied, so the sequence acts only through the believed link and the association differs at second order (IMM, ORD); under partial adoption the sequence decides which document is discounted, both ratings differ at order zero even for unrelated traits, and the association, $c$ times rescaling factors that now differ at order zero, differs at first order with coefficient $(1-\omega)H/Z$ (LAD). Propositions ORD and LAD must be applied (5.2, 6.3) before the references resolve.} \\
\multicolumn{3}{|p{\fullw}|}{\textbf{Evidence.} PropIMM.lean; PropORD.lean; Ladder.lean (\texttt{ladder\_gap}, \texttt{ladder\_gap\_mA1}, \texttt{ladder\_assoc\_coeff}, \texttt{ladder\_assoc\_coeff\_witness}); LemmaSEP.lean (general N, any attribute-local rescaling); sympy/verify\_ladder.py (53/53); sympy/verify\_interior\_omega.py rows 3 and 4.} \\ \hline
\multicolumn{3}{|p{\fullw}|}{\textbf{Grounds.}\par \textbf{Computation.} \textit{sympy/verify\_interior\_omega.py, check (3)} \newline Between-sequence effect on $\mathrm{assoc}$ with the second cue
adopted with weight $\delta$ (the manuscript's $\omega$): the $c^0$ term is $0$
for every $\delta$; the $c^1$ coefficient is $(1-\delta)\,G$ with $G$ affine in
$\delta$, $G\neq0$ at the generic prior for $\delta=1/2$ and for $\delta=0$; at
$\delta=1$ the factor $(1-\delta)$ makes the effect second order in $c$. \newline \textit{Under partial adoption the cross-product association carries the
sequence at first order in $c$.}} \\
\multicolumn{3}{|p{\fullw}|}{\textbf{Computation.} \textit{sympy/verify\_interior\_omega.py, check (4)} \newline Pooled belief $\lambda P_{AB} + (1-\lambda)P_{BA}$ at $c=0$:
$\mathrm{assoc}(\bar P) - \mathrm{assoc}(P^B)
= -\lambda(1-\lambda)(1-\delta)^2(\alpha-q_0)(\beta-r_0)$, nonzero for interior
$\lambda,\delta$ whenever both cues differ from the prior marginals, and zero
at $\delta=1$. \newline \textit{The pooled association is distorted already at $c=0$ unless the
later cue is adopted in full.}} \\
\multicolumn{3}{|p{\fullw}|}{\textbf{Computation.} \textit{sympy/verify\_interior\_omega.py, check (2)} \newline Each damped step multiplies the table it meets by a factor constant
along rows (resp.\ columns), so both routes are separable reweightings of the
prior and $\mathrm{OR}(P_{AB}) = \mathrm{OR}(P_{BA}) = \mathrm{OR}(P)$
for every $\delta$ (witnessed at $\delta=1/3$, $c=1/40$, generic prior). \newline \textit{Only the odds ratio is identical across sequences for every
adoption weight.}} \\
\multicolumn{3}{|p{\fullw}|}{\textbf{Theorem.} \textit{LemmaSEP.lean, isSeparable\_applySteps and sep\_rescales\_association} \newline For $N$ binary attributes, any finite sequence of attribute-local
updates gives $P'(x) = P(x)\prod_a g_a(x_a)$; and for $i\neq j$ such an $F$
satisfies $F_{11}F_{00} - F_{10}F_{01} = g_i(1)g_i(0)g_j(1)g_j(0)\,
\big(\prod_{a\neq i,j} g_a(x_a)\big)^2\,(P_{11}P_{00} - P_{10}P_{01})$. \newline \textit{Lemma SEP: a route rescales the association and cannot shift it,
whence the odds ratio is invariant; the cross-product itself is not.}} \\
\multicolumn{3}{|p{\fullw}|}{\textbf{Theorem.} \textit{PropDEC.lean, oddsRatio\_gap\_eq\_zero} \newline For $\alpha,\beta,q_0,r_0\notin\{0,1\}$, nonzero intermediate
$B$-marginals, nonzero off-diagonal prior cells and nonzero benchmark
normaliser and off-diagonal cells:
$\mathrm{OR}(P^J_{AB}) - \mathrm{OR}(P^B) = 0$, identically in $c$. \newline \textit{The odds-ratio form of the panel sentence is exact; the
cross-product form is protected only at $\omega=1$ (Proposition DEC).}} \\ \hline
\multicolumn{3}{|l|}{} \\ \hline
\entryhead{1.5}{Introduction, new paragraph after the hiring-panel paragraph (was 1.C)\quad{\footnotesize\textit{(formerly E.1)}}}
 & \textit{New paragraph(s) after the paragraph containing} ``with the difference set by arrival sequence.'' & Which impression prevails matters beyond the single judgement. The impression-formation literature disagrees on the direction. \citet{Asch1946} reports that the earlier terms set the direction in which later ones are read, and the review of \citet{HogarthEinhorn1992} finds primacy, recency or no effect depending on the characteristics of the task. Under either direction the sequence in which cues arrive is not neutral, and it need not be assigned by chance. A candidate who comes through a referral is met first through the letter and then through the credential, while a candidate who applies unsolicited is met first through the credential. If referral is more common in one group than in another, the two groups differ in the share $\lambda$ of evaluators who read the credential first. Two groups presenting the same evidence to evaluators with the same priors and the same preferences then receive different mean beliefs, and the difference is the sequence effect multiplied by the difference in their shares. A gap of that kind is not statistical discrimination in the sense of \citet{Phelps1972} and \citet{Arrow1973}, which rests on different prior beliefs about the groups, and it is none of the three sources that \citet{Bohren2019} distinguish, correct beliefs, biased beliefs and preferences, since none of these differs between the groups here. Which group it favours depends on whether the later or the earlier impression prevails, which Proposition~\ref{prop:ADJ} reads from the marginals. The gap also inherits the classification of this paper. It appears in the marginal probabilities and in the share of decisions changed, and an audit of the believed association does not register it. \\ \hline
\multicolumn{3}{|p{\fullw}|}{\textbf{Why.} Plan 1.C with audit P4 (``prior beliefs about the groups''), D13 (BIR's three sources named rather than ``none of the sources'', since their impartial type discriminates with correct beliefs and no animus), A24 (Asch's ``early terms dominate'' softened to his own ``direction'' account) and H6 (the primacy/recency/no-effect finding is HE's review of other studies) applied. Presupposes 5.2 (Proposition ADJ) and the Bohren2019 bib entry. Goes after 1.4's sentence.} \\
\multicolumn{3}{|p{\fullw}|}{\textbf{Evidence.} The identity mean(lambda) - mean(lambda') = (lambda - lambda')(PJ\_AB - PJ\_BA) is exact, verify\_ORD.py step 9; Asch.lean; HogarthEinhorn.lean; BohrenImasRosenberg.lean.} \\ \hline
\multicolumn{3}{|p{\fullw}|}{\textbf{Grounds.}\par \textbf{Computation.} \textit{sympy/verify\_ORD.py, step (9), with steps (1), (3)--(5)} \newline With $\Pbar_\lambda=\lambda\,\PJ_{AB}+(1-\lambda)\,\PJ_{BA}$, the identity
$\Pbar_{\lambda}-\Pbar_{\lambda'}=(\lambda-\lambda')\,(\PJ_{AB}-\PJ_{BA})$ holds entrywise
and exactly, at every order in $c$ (symbolic in $\alpha,\beta,c,q_0,r_0,\lambda,\lambda'$).
Since $\PJ_{AB}-\PJ_{BA}=c\,(\kappa R_1-\kappa' R_2)+\bigO(c^2)$, the group gap on the
$A$-marginal is $(\lambda-\lambda')\,\kappa'\,c+\bigO(c^2)$ with
$\kappa'=(\beta-r_0)q_0(1-q_0)/Z$, while on the believed association it is $\bigO(c^2)$
(checked at $\lambda=2/5$ against $\lambda'=3/4$). \newline \textit{Two groups with the same priors, preferences and evidence differ by the
sequence effect multiplied by the difference in their shares, and the gap inherits
Proposition ORD's classification (marginals first order, association second order).}} \\
\multicolumn{3}{|p{\fullw}|}{\textbf{Quote.} \textit{Asch 1946, pp.~271--272 (Experiment VI)} \newline ``the first terms set up in most subjects a direction which then exerts a
continuous effect on the latter terms'' (p.~271); ``a broad, uncrystallized but directed
impression is born. The next characteristic comes not as a separate item, but is related
to the established direction. Quickly the view formed acquires a certain stability, so
that later characteristics are fitted\ldots to the given direction'' (p.~272). \newline \textit{Asch's account is one of direction rather than dominance, and he adds that it is
``not the sheer temporal position of the item which is important'' (p.~272), which is why
the entry no longer says that early terms dominate (audit A24).}} \\
\multicolumn{3}{|p{\fullw}|}{\textbf{Quote.} \textit{Hogarth and Einhorn 1992, abstract, T-p.~1 (transcription)} \newline ``Much literature attests to the existence of order effects in the updating of
beliefs. However, under what conditions do primacy, recency, or no order effects occur?
This paper presents a theory of belief updating that explicitly accounts for order-effect
phenomena as arising from the interaction of information-processing strategies and task
characteristics.'' \newline \textit{The three outcomes are their review's classification of earlier studies (Table 1,
T-p.~4) and their model's predictions (Table 2); their own five experiments found recency
or no effect, so the entry attributes the finding to the review (audit H6).}} \\
\multicolumn{3}{|p{\fullw}|}{\textbf{Quote.} \textit{Bohren, Imas and Rosenberg, January 2019 working paper, pp.~2, 13--14, 19--20} \newline ``we allow for three potential sources: (i) belief-based with correct beliefs,
(ii) belief-based with incorrect, biased beliefs, and (iii) preference-based'' (p.~2).
Of the first, ``discrimination is due to imperfect information (Phelps 1972)'' or ``is an
equilibrium effect (Arrow 1973)'' (pp.~13--14), and ``evaluators hold prior beliefs about
workers' abilities that differ by group identity'' (p.~14). The impartial type ``has no
belief-based partiality'' (p.~19), yet ``this type discriminates against males in the
second period'' (p.~20). \newline \textit{The three sources are named rather than treated as exhaustive, since BIR's own
impartial type discriminates with correct beliefs and no animus (audit D13).}} \\
\multicolumn{3}{|p{\fullw}|}{\textbf{Quote.} \textit{Phelps 1972, p.~659; Arrow 1973, p.~26 (working paper)} \newline Phelps: ``the employer who seeks to maximize expected profit will discriminate
against blacks or women if he believes them to be less qualified, reliable, long-term,
etc. on the average than whites and men, respectively, and if the cost of gaining
information about the individual applicants is excessive.'' Arrow: ``The employer cannot
know of any given worker whether or not he is qualified; however, he does believe that the
probability that a random W worker is qualified is $p_W$ and that a random B worker is
qualified is $p_B$.'' \newline \textit{Both rest on prior beliefs about the groups, the wording the entry uses (audit P4);
Phelps's Case 2 and Further Case rest on variance and test reliability rather than on
mean, so ``different beliefs'' alone would be too broad.}} \\ \hline
\multicolumn{3}{|l|}{} \\ \hline
\entryhead{1.6}{Introduction, premise paragraph, from ``Read as a Bayes factor''\quad{\footnotesize\textit{(formerly C.6)}}}
 & Read as a Bayes factor, the credential carries a likelihood ratio stating how much more probable it is under competence than in its absence, and that ratio multiplies whatever belief it meets. A Bayes factor, however, requires the probability of the same credential for a candidate who is not competent, and an impression does not supply that counterfactual. As a mechanism that holds conditionals fixed, Jeffrey updating is the unique coherent revision for credence delivered on the cue's partition and is equivalent to making the minimal change of the prior consistent with the delivered marginal \citep{DiaconisZabell1982} & Read as a Bayes factor, the credential carries the ratio of new to old odds on competence, which multiplies whatever belief it meets. The two readings agree on a single cue (Proposition~\ref{prop:IMM}) and differ in what stays fixed when the cue meets another prior. The panelist's own impression is naturally a credence, whereas a factor suits evidence reported by someone else, whose report mixes the evidence with that person's prior \citep[ch.~3]{Jeffrey2004}. Jeffrey updating holds fixed the conditionals given the cue's partition, and its posterior is the unique belief with the delivered marginal closest to the prior in Kullback--Leibler divergence \citep[Theorem~5.1]{DiaconisZabell1982} \\ \hline
\multicolumn{3}{|p{\fullw}|}{\textbf{Why.} (i) ``A Bayes factor requires the probability of the same credential for a candidate who is not competent'' is contradicted by Jeffrey's own definition, new odds over old odds, which needs no likelihood (2002 draft of Jeffrey 2004, ch. 3; audit J2). Jeffrey draws the line between the readings by provenance instead (own experience gives credences, others' reports should be converted to factors), which supports the paper's modelling choice. (ii) ``the unique coherent revision'' is not in Diaconis-Zabell; what they prove is the unique Kullback-Leibler (and Hellinger) minimiser (M22). ``equivalent to making the minimal change'' without naming the distance is also loose (not unique in variation distance). The footnote after the citation is unchanged. Needs the Jeffrey2004 bib entry (B.2); the Drive copy is the 2002 draft, so cite the chapter, not a page.} \\
\multicolumn{3}{|p{\fullw}|}{\textbf{Evidence.} verify\_jeffrey.md J2; DiaconisZabell.lean (\texttt{thm51\_KL\_eq\_iff}, \texttt{jcond\_iff\_jeffrey}); PropIMM.lean (\texttt{propIMM\_single\_cue}).} \\ \hline
\multicolumn{3}{|p{\fullw}|}{\textbf{Grounds.}\par \textbf{Quote.} \textit{Jeffrey (2002 draft of Jeffrey 2004), ch.~3, \S3.3, p.~61, eq.~(4)} \newline ``$\beta(D_i : D_1) = \dfrac{new\,D_i}{new(D_1)} \Big/ \dfrac{old(D_i)}{old(D_1)}
= \dfrac{\pi(D_i)}{\pi(D_1)}$ \quad ($D_i : D_1$ Bayes factor). This is your Bayes factor
for $D_i$ against $D_1$; \ldots it is what remains of your new odds when the old odds
have been factored out.'' And: ``probability factors are a little easier to compute
than Bayes factors, starting from old and new probabilities.'' \newline \textit{Jeffrey's Bayes factor is new odds over old odds, computed from the prior and the delivered credence; no likelihood under a counterfactual enters, so ``requires the probability of the same credential for a candidate who is not competent'' is false of it.}} \\
\multicolumn{3}{|p{\fullw}|}{\textbf{Quote.} \textit{Jeffrey (2002 draft), ch.~3, \S3.2 p.~59 and \S3.3 p.~59} \newline ``But the results $new'(D_i)$ of that interaction are \textit{data} for the
kinematical formula for updating $old'(H)$; the formula itself does not compute those
results.'' ``You are unlikely to simply adopt such an observer's updated probabilities
as your own, for they are necessarily a confusion of what the other person has gathered
from the observation itself \ldots with that person's prior judgmental state, for which
you may prefer to substitute your own.'' \newline \textit{Jeffrey draws the line by provenance: one's own experience yields credences taken as data, another's report is to be converted to factors (\S3.3), which is what the corrected sentence says.}} \\
\multicolumn{3}{|p{\fullw}|}{\textbf{Quote.} \textit{Jeffrey (2002 draft), ch.~2, p.~42} \newline ``Assuming rigidity relative to $D$, the odds factor for a theory $T$ against an
alternative theory $S$ that is due to learning that $D$ is true will be the left-hand side
of the following equation, the right-hand side of which is called `the likelihood ratio':
Bayes Factor = Likelihood Ratio: $\dfrac{old(T|D)/old(S|D)}{old(T)/old(S)} =
\dfrac{old(D|T)}{old(D|S)}$'' \newline \textit{The Bayes factor equals a likelihood ratio only when $D$ is learned for certain under rigidity; the counterfactual-likelihood reading the manuscript objects to is this special case.}} \\
\multicolumn{3}{|p{\fullw}|}{\textbf{Theorem.} \textit{DiaconisZabell.lean, thm51\_KL\_eq\_iff (with thm51\_KL\_le); D-Z Theorem 5.1, (5.6)} \newline Let $\Omega$ be finite, $P(\omega)>0$ for all $\omega$, $e:\Omega\to I$ surjective
onto a finite $I$ with cells $E_i=e^{-1}(i)$, and $Q\ge0$ with $Q(E_i)=p_i$ for every $i$.
Then $\mathrm{KL}(Q\,\|\,P)=\sum_\omega Q(\omega)\log\frac{Q(\omega)}{P(\omega)}
\;\ge\;\sum_i p_i\log\frac{p_i}{P(E_i)}$, with equality if and only if
$Q(\omega)=p_{e(\omega)}\,P(\omega)/P(E_{e(\omega)})$, Jeffrey's update. \newline \textit{What D-Z prove is that Jeffrey's posterior is the unique Kullback--Leibler minimiser among beliefs with the delivered marginal; ``the unique coherent revision'' is not in the paper.}} \\
\multicolumn{3}{|p{\fullw}|}{\textbf{Theorem.} \textit{DiaconisZabell.lean, remark\_a\_tv\_not\_unique; D-Z Remark (a), p.~828} \newline On $\Omega=\{0,1,2\}$ with cells $\{0,1\}$, $\{2\}$, prior $P=(1/4,1/4,1/2)$ and new
cell probabilities $(3/4,1/4)$: Jeffrey's update is $J=(3/8,3/8,1/4)$; $Q=(1/2,1/4,1/4)$
has the same cell probabilities, $Q\neq J$, and $\|Q-P\|=\|J-P\|=\tfrac12\sum_i|p_i-P(E_i)|=1/4$,
the minimum of (5.4). D-Z: ``Although the probability measure given by Jeffrey's rule
minimizes the variation distance, it does not do so uniquely''. \newline \textit{``Equivalent to making the minimal change'' without naming the distance is loose: in variation distance the minimiser is not unique.}} \\
\multicolumn{3}{|p{\fullw}|}{\textbf{Computation.} \textit{sympy/verify\_soft\_vs\_hard.py (15/15 checks pass)} \newline Worked example $\alpha=\beta=1/2$, $c=1/20$, $q_0=1/5$, $r_0=7/10$. The letter's
implied factor (new odds over old odds on $B$) is $\frac{7/3}{1}=7/3$ read first, against
$\beta=1/2$, and $\frac{7/3}{11/14}=98/33$ read second, against the post-credential
$B$-marginal $11/25$. Hard cues with fixed ratios $1/4$ and $7/3$ commute cell by cell and
reproduce $P^{B}$, with $A$-marginal $27/119$ in either sequence; the soft sequences end at
$18/77$ (credential first) and $1/5$ (letter first). \newline \textit{The two readings differ in what is held fixed when the cue meets a moved marginal, not in whether a counterfactual likelihood exists.}} \\
\multicolumn{3}{|p{\fullw}|}{\textbf{Theorem.} \textit{JeffreyOrder/PropIMM.lean, propIMM\_single\_cue (with bayesAW\_total)} \newline For $\alpha\neq0$, $1-\alpha\neq0$ and every $\beta$, $c$, $q_0$:
$\mathrm{bayesA}(P(\alpha,\beta,c),\alpha,q_0)=\mathrm{jeffreyA}(P(\alpha,\beta,c),q_0)$,
i.e. reweighting $P$ by $q_0/\alpha$ on row $A{=}0$ and $(1-q_0)/(1-\alpha)$ on row
$A{=}1$ and normalising equals the Jeffrey step to $A$-marginal $q_0$, identically; the
normaliser is $1$. \newline \textit{On a single cue the two readings agree, as the corrected sentence says (Proposition IMM).}} \\ \hline
\multicolumn{3}{|l|}{} \\ \hline
\entryhead{1.7}{Introduction, terminology paragraph, Domotor sentence and two typos\quad{\footnotesize\textit{(formerly C.7)}}}
1/2 & The sequence-dependence of marginals follows easily -- since the likelihood a delivered credence implies must be read against the marginal in force when the cue arrives \citep{Domotor1980}. & The sequence dependence of the marginals follows directly, since a delivered credence is taken against the marginal in force when the cue arrives, which the first cue moves when $c\neq0$ (Proposition~\ref{prop:DIV}). \\ \hline
2/2 & To avoid a clash in terminology we use ``order'' to describe proximity with a sequence-free Bayesian bechmark and ``sequence`` to denote the sequence in which evidence is read. & To avoid a clash in terminology we use ``order'' to describe proximity with a sequence-free Bayesian benchmark and ``sequence'' to denote the sequence in which evidence is read. \\ \hline
\multicolumn{3}{|p{\fullw}|}{\textbf{Why.} Domotor never mentions likelihoods or marginals; he supports only bare non-commutativity, and treats it as a defect of Jeffrey machines (p. 395), the opposite lean (audit M1). The mechanism is the paper's own (Proposition DIV). If a Domotor citation is wanted, it fits the literature paragraph as a source of the commutativity demand that Jeffrey (1983, pp. 182-183) rejects. Typos: ``bechmark'', \texttt{}sequence\texttt{} closed with backquotes.} \\
\multicolumn{3}{|p{\fullw}|}{\textbf{Evidence.} Domotor.lean (\texttt{field\_eq\_jeffrey\_embed}, \texttt{embed\_depends\_on\_state}); Cripps.lean (\texttt{composite\_AB\_eq\_bayes}: the B-marginal moves by c(q0-alpha)/(alpha(1-alpha))).} \\ \hline
\multicolumn{3}{|p{\fullw}|}{\textbf{Grounds.}\par \textbf{Quote.} \textit{Domotor (1980), p.~395} \newline ``Along with Field (1978) we may argue that Jeffrey machines are inadequate because,
among other things, the commutativity of evidence application:
$[P_{(U,p)}]_{(V,q)} = [P_{(V,q)}]_{(U,p)}$ fails in general.'' \newline \textit{Domotor supports bare non-commutativity and counts it against Jeffrey machines; the manuscript cites him for a mechanism and for the opposite lean.}} \\
\multicolumn{3}{|p{\fullw}|}{\textbf{Quote.} \textit{Domotor (1980), p.~397} \newline ``A moment's thought shows that Field's input space is embeddable into Jeffrey's
input space by the map: $h_P : \mathbf{F}_X \to \mathbf{E}_X$, defined by
$h_P(U,\alpha) = (U,p)$, where: $p(A) = e^{\alpha_A} P(A) : \alpha_X$.'' \ldots ``So,
whatever advantage may have been gained in commutativity is lost in probabilistic
independence.'' \newline \textit{The passage nearest a mechanism: a fixed factor $\alpha$ yields a $p$ that depends on the state $P$, used to argue that Field's space is smaller, not to explain non-commutativity. Neither ``likelihood'' nor ``marginal'' occurs on the rendered pp.~394--397, the pages around the quoted passages.}} \\
\multicolumn{3}{|p{\fullw}|}{\textbf{Theorem.} \textit{Domotor.lean, field\_eq\_jeffrey\_embed; embed\_depends\_on\_state} \newline For $P\ge0$ on finite $\Omega$, a partition $u:\Omega\to I$ and $\alpha:I\to\mathbb{R}$,
put $\alpha_X=\sum_y e^{\alpha_{u(y)}}P(y)$ and $h_P(\alpha)_i=e^{\alpha_i}P(E_i)/\alpha_X$.
Then Field's conditional $x\mapsto e^{\alpha_{u(x)}}P(x)/\alpha_X$ equals Jeffrey's rule on
$(u,h_P(\alpha))$. The image depends on the state: $\alpha=(\log2,0)$ on two atoms gives
$p=(2/3,1/3)$ at $P=(1/2,1/2)$ but $p=(2/5,3/5)$ at $P=(1/4,3/4)$. \newline \textit{Read in reverse, a fixed delivered credence corresponds to a factor that depends on the belief it meets; Domotor states the embedding but draws neither reading.}} \\
\multicolumn{3}{|p{\fullw}|}{\textbf{Theorem.} \textit{Cripps.lean, composite\_AB\_eq\_bayes} \newline Let $\mu$ on $A\times B$ be strictly positive with $\sum\mu=1$, $q$ on $A$ strictly
positive with $\sum_a q_a=1$, and $r$ on $B$ with $\sum_b r_b=1$. Then
$J_B(J_A(\mu,q),r)=\mathrm{bayes}\big(\mu,\ x\mapsto
\tfrac{q_{x_1}}{\mu(A{=}x_1)}\cdot\tfrac{r_{x_2}}{(J_A(\mu,q))(B{=}x_2)}\big)$, where
$\mathrm{bayes}(\mu,\ell)(\theta)=\mu(\theta)\ell(\theta)/\sum_{\theta'}\mu(\theta')\ell(\theta')$
and $J_A(\mu,q)(x)=q_{x_1}\mu(x)/\mu(A{=}x_1)$. In the other sequence the $B$-likelihood is
$r_{x_2}/\mu(B{=}x_2)$. \newline \textit{The second cue's likelihood is matched to the intermediate belief, which the first cue has moved; this is the corrected sentence's mechanism, and it is the paper's own.}} \\
\multicolumn{3}{|p{\fullw}|}{\textbf{Computation.} \textit{sympy, recomputed symbolically (the instance is check 3 of verify\_soft\_vs\_hard.py)} \newline After the $A$-step to $q_0$ from $P(\alpha,\beta,c)$, the $B$-marginal is
$\beta + c\,\dfrac{q_0-\alpha}{\alpha(1-\alpha)}$, so the first cue moves it whenever
$c\neq0$ and $q_0\neq\alpha$. At $\alpha=\beta=1/2$, $c=1/20$, $q_0=1/5$ it is $11/25$. \newline \textit{The delivered credence $r_0$ is then read against $11/25$ instead of $1/2$, which is the whole of the sequence effect.}} \\
\multicolumn{3}{|p{\fullw}|}{\textbf{Quote.} \textit{PAPER\_B\_MANUSCRIPT.tex, Proposition DIV (micro divergence), (i)} \newline Let $Z:=\alpha\beta(1-\alpha)(1-\beta)>0$. ``At $c=0$ both reading sequences coincide
with the benchmark, $P^{J}_{AB}=P^{J}_{BA}=P^{B}=q\otimes r$.'' (i) For $c\neq0$,
$P^{J}_{AB}-P^{B}=c\,\kappa R_1+O(c^{2})$, $R_1=q\otimes(1,-1)$,
$\kappa=\frac{(\alpha-q_0)\,r_0(1-r_0)}{Z}$, ``with $\kappa=0\iff q_0=\alpha$ or
$r_0\in\{0,1\}$''. \newline \textit{The manuscript's own statement of the mechanism the sentence should cite.}} \\ \hline
\multicolumn{3}{|l|}{} \\ \hline
\entryhead{1.8}{Introduction, key-finding paragraph, last two sentences\quad{\footnotesize\textit{(formerly C.8)}}}
 & The implication for evaluation is that and audit's answer depends on which statistic is aggregated rather than on how many evaluators are averaged over. One consequence, is that the instruments most naturally trusted for detecting stereotype---belief audits asking how attributes are thought to go together---may not even be detectable. & The implication for evaluation is that an audit's answer depends on which statistic is aggregated rather than on how many evaluators are averaged over. One consequence is that belief audits asking how attributes are thought to go together, the instruments most naturally trusted for detecting stereotype, cannot detect the sequence effect at first order. \\ \hline
\multicolumn{3}{|p{\fullw}|}{\textbf{Why.} Typos (``and audit's'', ``One consequence, is''), and ``the instruments ... may not even be detectable'' says the instruments are undetectable; it is the sequence effect that belief audits cannot detect at first order.} \\
\multicolumn{3}{|p{\fullw}|}{\textbf{Evidence.} Propositions DEC, PRO.} \\ \hline
\multicolumn{3}{|p{\fullw}|}{\textbf{Grounds.}\par \textbf{Theorem.} \textit{PropDEC.lean, assoc\_PJab; PropPRO.lean, lemmaASC\_first\_order\_vanishes} \newline $\mathrm{assoc}(P^J_{AB}) = \frac{q_0}{\alpha}\frac{1-q_0}{1-\alpha}
\frac{r_0}{m_0}\frac{1-r_0}{m_1}\,c$ with $m_j$ the $B$-marginal after the
$A$-step; and along either route direction
$\mathrm{assoc}(q\otimes r + t\kappa R_\sigma) = \mathrm{assoc}(q\otimes r)
+ t^2\,\mathrm{assoc}(\kappa R_\sigma)$. \newline \textit{Proposition DEC: the believed association departs from the
benchmark only at second order in $c$.}} \\
\multicolumn{3}{|p{\fullw}|}{\textbf{Theorem.} \textit{PropPRO.lean, propPRO\_protection and margA\_not\_in\_span} \newline $\langle xJ + y\,\nabla\mathrm{assoc},\,\lambda\kappa R_1 + (1-\lambda)\kappa' R_2\rangle = 0$
for all $x,y,\alpha,\beta,\lambda$; whereas for $0<r_0<1$ the marginal read-out
$(0,0,1,1)$ is not in $\mathrm{span}\{J,\nabla\mathrm{assoc}\}$. \newline \textit{Proposition PRO: which statistic is aggregated decides whether the
sequence effect is first or second order.}} \\
\multicolumn{3}{|p{\fullw}|}{\textbf{Theorem.} \textit{PropPRO.lean, propPRO\_uniqueness} \newline For $\lambda\in(0,1)$, $q_0,r_0\in(0,1)$, $\beta\neq r_0$: if the
aggregate first-order coefficient
$\lambda(\alpha-q_0)r_0(1-r_0)x + (1-\lambda)(\beta-r_0)q_0(1-q_0)y$ vanishes at
two priors $\alpha_1\neq\alpha_2$, then $x=y=0$. \newline \textit{The coefficient is a property of the population's belief, not of a
sample; averaging over more evaluators does not change it.}} \\ \hline
\multicolumn{3}{|l|}{} \\ \hline
\entryhead{1.9}{Introduction, roadmap, the Section 6 sentence\quad{\footnotesize\textit{(formerly E.15r)}}}
 & Section~\ref{sec:scope} states the scope of results and limitations that would break the results. & Section~\ref{sec:scope} states the scope of the results and the limitations that would break them, shows that the second-order results require full adoption of the later cue, and gives a test of that condition from the ratings. \\ \hline
\multicolumn{3}{|p{\fullw}|}{\textbf{Why.} The roadmap's Section 6 sentence names what that section now does: the scope of the results, that the second-order results require full adoption of the later cue, and the test of that condition (6.3).} \\
\multicolumn{3}{|p{\fullw}|}{\textbf{Evidence.} Proposition ADJ (6.3).} \\ \hline
\multicolumn{3}{|l|}{} \\ \hline
\multicolumn{3}{|l|}{\rule{0pt}{3.2ex}\Large\textbf{2\quad Setup}} \\ \hline\hline
\entryhead{2.6}{Setup 2.3, Assumption 2 (was 2.C)\quad{\footnotesize\textit{(formerly E.4)}}}
 & The cues are bundled, correlated, and soft, each delivering a credence on an attribute's partition, not a decisive cue. & The cues are bundled, correlated, and soft, each delivering a credence on an attribute's partition, not a decisive cue, with $0<q_0<1$ and $0<r_0<1$. \\ \hline
\multicolumn{3}{|p{\fullw}|}{\textbf{Why.} Plan 2.C, with the colon replaced by a clause.} \\
\multicolumn{3}{|p{\fullw}|}{\textbf{Evidence.} Notation only.} \\ \hline
\multicolumn{3}{|p{\fullw}|}{\textbf{Grounds.}\par \textbf{Theorem.} \textit{PropDIV.lean, kappa\_eq\_zero\_iff and kappa'\_eq\_zero\_iff} \newline With $Z=\alpha\beta(1-\alpha)(1-\beta)\neq0$,
$\kappa=(\alpha-q_0)\,r_0(1-r_0)/Z$ and $\kappa'=(\beta-r_0)\,q_0(1-q_0)/Z$:
$\kappa=0 \iff \alpha=q_0 \ \text{or}\ r_0=0 \ \text{or}\ r_0=1$, and
$\kappa'=0 \iff \beta=r_0 \ \text{or}\ q_0=0 \ \text{or}\ q_0=1$. \newline \textit{A cue with $r_0\in\{0,1\}$ is decisive and kills the leading gap, so ``soft'' means
$0<q_0<1$ and $0<r_0<1$, the interior credences Assumption 2 now states.}} \\
\multicolumn{3}{|p{\fullw}|}{\textbf{Theorem.} \textit{PropPRO.lean, propPRO\_uniqueness (hypotheses)} \newline Proposition PRO (ii) is proved under $0<\lambda<1$, $0<q_0<1$, $0<r_0<1$,
$\beta\neq r_0$ and two priors $\alpha_1\neq\alpha_2$; the interior bounds on $q_0,r_0$
are what make the factors $r_0(1-r_0)$ and $q_0(1-q_0)$ nonzero in the leading
coefficient $\lambda(\alpha-q_0)r_0(1-r_0)\,x+(1-\lambda)(\beta-r_0)q_0(1-q_0)\,y$. \newline \textit{The same interior condition on the credences is a hypothesis of the paper's
central characterisation.}} \\ \hline
\multicolumn{3}{|l|}{} \\ \hline
\multicolumn{3}{|l|}{\rule{0pt}{3.2ex}\Large\textbf{3\quad Related literature}} \\ \hline\hline
\entryhead{3.2}{Related literature, the paragraph from BHW/Banerjee to Becker and Good-Mittal\quad{\footnotesize\textit{(formerly C.11)}}}
1/5 & In \citet{Banerjee1992}, the sequence persists because a coarse public action fails to be a sufficient statistic for private information. In \citet{BHW1992}, on the other hand, a cascade is an action that conveys no private signal. & In both, the equilibrium action is not a sufficient statistic for private information, being binary in \citet{BHW1992} and a non-invertible choice rule in \citet{Banerjee1992}, so once a cascade or herd forms, later actions convey no private signal. \\ \hline
2/5 & One similarity with the literature, however, is that the sequence effect in \citep{BHW1992, Banerjee1992} lives on the analogue of the marginals---which remain so in the current paper as well. & One analogy remains. Those models have a single state, so their sequence effect falls on a belief about one variable, the counterpart of a marginal, and here too the marginals carry the sequence effect. \\ \hline
3/5 & A related normative literature evaluates group belief formation by whether aggregation commutes with updating---a design criterion for the pooling rule \citep{Dietrich2021}. While geometric pooling \citep[Def.~1]{Dietrich2021} satisfies the criterion, the population mean studied in the current paper (i.e.\ the average posterior across evaluators who differ only in the sequence they meet the two cues) is the outcome of linear pooling and generically fails the criterion. As further sections elaborate, this is consistent with the first-order divergence of every statistic outside the protected class of Proposition~\ref{prop:PRO}. & A related normative literature evaluates group belief formation by whether aggregation commutes with updating, a criterion that geometric pooling meets and linear pooling generally does not \citep[Theorem~2]{Dietrich2021}. The criterion does not discriminate here. Evaluators who share a prior and differ only in reading sequence do not hold conditionalisations of that prior on a common event, and the linear and geometric averages of their posteriors differ only at $\bigO(c^{2})$. The first-order divergence of every statistic outside the protected class of Proposition~\ref{prop:PRO} is therefore a property of Jeffrey updating, not of linear pooling. \\ \hline
4/5 & \textbf{Several alternatives to Bayesian updating have been used in the decision theory literature to explain whether sequential conditioning is sequence-independent or not}. For example, \citet{Epstein2006} makes the updating rule subjective, and \citet{Ortoleva2012} axiomatises departures triggered by unexpected news. \citet{Cripps2021} shows that symmetry and divisibility jointly force sequence-independence, so the correlated two-cue composite of Proposition~\ref{prop:DIV} fails divisibility---and, of the four axioms, only divisibility.\footnote{Cripps's axioms govern rules on fixed-likelihood experiments; the statement here is behavioural---the map from prior and delivered credences to final belief, held against the sequence-invariance that his Symmetry and Divisibility axioms jointly force.} & Decision theory has axiomatised several alternatives to Bayesian updating, each for a single piece of news. \citet{Epstein2006} makes the updating rule, and not only the prior, subjective, and \citet{Ortoleva2012} axiomatises departures triggered by unexpected news. \citet{Cripps2021} treats the sequence of signals, and his Symmetry and Divisibility axioms imply that two conditionally independent signals give the same posterior in either sequence. Read with each cue as its likelihood matched to the prior (Proposition~\ref{prop:IMM}), the composite of Proposition~\ref{prop:DIV} is a sequence of Bayes updates satisfying all four of his axioms, and it depends on the sequence because the second likelihood is matched to the intermediate belief, so the two sequences process different experiments. \\ \hline
5/5 & While the literature restores sequence-invariance for sequential Jeffrey updating by changing how successive inputs are pooled \citep{PettigrewWeisberg2025}, & While \citet{PettigrewWeisberg2025} restore sequence-invariance by pooling the prior with each new input multiplicatively before the Jeffrey step, \\ \hline
\multicolumn{3}{|p{\fullw}|}{\textbf{Why.} One entry per sentence group, in manuscript order. BHW/Banerjee: ``coarse'' belongs to BHW's binary action, not Banerjee's continuum; the two share one mechanism, so ``on the other hand'' goes; the marginals point is the paper's analogy (M2, M3). Dietrich: Def. 1 is preference aggregation, and with a common prior linear pooling passes the criterion wherever it applies (M4, M5). Epstein/Ortoleva/Cripps (bold sentence): the framing is unsupported for Epstein and Ortoleva; the Cripps claim is false, since under Proposition IMM's reading the composite satisfies all four axioms, and the footnote goes with it (M6-M9). Pettigrew-Weisberg pool the prior with each input, not successive inputs (M25). The rest of this paragraph is 3.4, since the Asch entry (3.3) falls in between.} \\
\multicolumn{3}{|p{\fullw}|}{\textbf{Evidence.} BHW.lean, Banerjee.lean, Dietrich.lean, Epstein.lean, Ortoleva.lean (via the 2024 survey), Cripps.lean (\texttt{order\_invariance}, \texttt{composite\_AB\_eq\_bayes}), PettigrewWeisberg.lean.} \\ \hline
\multicolumn{3}{|p{\fullw}|}{\textbf{Grounds.}\par \textbf{Quote.} \textit{Bikhchandani, Hirshleifer and Welch 1992, p. 1000} \newline ``DEFINITION. An informational cascade occurs if an individual's
action does not depend on his private information signal. If an individual $i$
is in a cascade, then his action conveys no information and individual $i+1$
draws the same inference from all previous actions.'' \newline \textit{(part 1, BHW) A cascade is defined by the action's silence about
the signal; the action is binary (adopt or reject, p. 996), which is where
``coarse'' belongs.}} \\
\multicolumn{3}{|p{\fullw}|}{\textbf{Quote.} \textit{Banerjee 1992, p. 809} \newline ``The key reason why we get a different result is that in our
model the choices made by agents are not always sufficient statistics for the
information they have. \ldots It is evident that, for a wide range of signals,
the agents in our model will always choose the same option; this lack of
invertibility is what causes the sufficient property to fail.'' \newline \textit{(part 1, Banerjee) The same mechanism as BHW, on a continuum of
options, so ``on the other hand'' has no contrast to mark.}} \\
\multicolumn{3}{|p{\fullw}|}{\textbf{Theorem.} \textit{BHW.lean, lik\_ratio (with cascade\_uninformative)} \newline For every $p$ and every history $h$ with
$d(h) = \#\mathrm{adopt} - \#\mathrm{reject}$,
$L_1(h)\,c_0(d) = L_0(h)\,c_1(d)$, so the public likelihood ratio
$L_1/L_0$ is a function of $d$ alone: $1$ at $d=0$, $p/(1-p)$ at $d=1$,
$p(1+p)/((1-p)(2-p))$ at $d\ge 2$, and the reciprocals at $d=-1$, $d\le-2$. \newline \textit{(part 1, BHW) Two adoptions carry less than two $H$ signals,
since $p(1+p)/((1-p)(2-p)) < (p/(1-p))^2 \iff 1 < 2p$; and the ratio is
constant once $|d|\ge 2$, so later actions convey no private signal.}} \\
\multicolumn{3}{|p{\fullw}|}{\textbf{Quote.} \textit{Bikhchandani, Hirshleifer and Welch 1992, p. 996} \newline ``The gain to adopting, $V$, is also the same for all
individuals and is either zero or one, with equal prior probability 1/2.'' \newline \textit{(part 2, marginals analogy) One binary state, so the sequence
effect in BHW falls on a belief about one variable, the counterpart of a
marginal.}} \\
\multicolumn{3}{|p{\fullw}|}{\textbf{Quote.} \textit{Banerjee 1992, p. 802} \newline ``Let us assume that there is a unique $i^{*}$ such that
$z(i) = 0$ for all $i \neq i^{*}$ and $z(i^{*}) = z$, where $z > 0$. \ldots
Of course, everybody, given these payoffs, would want to invest in $i^{*}$.
The trouble is no one knows which one it is.'' \newline \textit{(part 2, marginals analogy) One unknown, the identity of
$i^{*}$; there is no second attribute and no association for the sequence
effect to fall on, which is the analogy the AFTER keeps.}} \\
\multicolumn{3}{|p{\fullw}|}{\textbf{Quote.} \textit{Dietrich 2021, Definition 1, p. 7} \newline ``Definition 1 An EU preference aggregation rule \ldots is
linear-geometric if there exist individual weights \ldots such that for each
preference profile \ldots the group preference relation \ldots has \ldots
probability function \ldots given on states by \ldots up to a multiplicative
constant.'' \newline \textit{(part 3, Dietrich) Definition 1 defines a rule on preference
profiles, so it is preference aggregation; the belief-pooling statement is
Theorem 2, which the AFTER cites.}} \\
\multicolumn{3}{|p{\fullw}|}{\textbf{Quote.} \textit{Dietrich 2021, Theorem 2, p. 11} \newline ``Theorem 2 A belief aggregation rule (on the domain of coherent
belief profiles) is dynamically rational, unanimity-preserving, and continuous
if and only if it is geometric.'' \newline \textit{(part 3, Dietrich) The criterion is Dynamic Rationality, defined
on p. 16 as ``When information is learnt by everyone, the new group beliefs
equal the old ones conditional on the information''.}} \\
\multicolumn{3}{|p{\fullw}|}{\textbf{Theorem.} \textit{Dietrich.lean, linPool\_dynRational\_on\_common\_prior} \newline Let $w_i \ge 0$ with $\sum_i w_i = 1$, let $p_i = \pi$ for every
$i$, and let $E$ be any event. Then
$\sum_i w_i\, p_i(\cdot \mid E) = \big(\sum_i w_i\, p_i\big)(\cdot \mid E)$. \newline \textit{(part 3, Dietrich) With a common prior the premise of Dynamic
Rationality forces unanimity, and linear pooling passes wherever the criterion
applies; Paper B's evaluators do not condition on a common event, so the
criterion is silent on them.}} \\
\multicolumn{3}{|p{\fullw}|}{\textbf{Quote.} \textit{Epstein 2006, abstract, p. 413} \newline ``The main result is a representation theorem that generalizes
(the dynamic version of) Anscombe--Aumann's theorem so that both the prior and
the way in which it is updated are subjective.'' \newline \textit{(part 4, Epstein) What Epstein makes subjective is the
updating rule, not only the prior; nothing in the abstract concerns the
sequence of several updates.}} \\
\multicolumn{3}{|p{\fullw}|}{\textbf{Quote.} \textit{Epstein 2006, p. 416} \newline ``There are three periods---an ex ante stage 0, an interim period
1 when a signal $s_1 \in S_1$ is realized, and period 2 when the remaining
uncertainty is resolved through realization of some $s_2 \in S_2$.'' \newline \textit{(part 4, Epstein) One interim signal, so the model has no
sequence of updates whose invariance could be at issue.}} \\
\multicolumn{3}{|p{\fullw}|}{\textbf{Quote.} \textit{Ortoleva 2024, p. 546 and footnote 2} \newline ``we do not discuss updating \ldots when subjects receive more
general types of information, such as vague, ambiguous, or imprecise
information or information of the form `Event A is more likely than event
B';'' with footnote 2: ``These cases have been extensively studied in several
disciplines and include classical approaches such as Jeffrey's rule or the AGM
theory of belief revisions''. \newline \textit{(part 4, Ortoleva) Jeffrey's rule is outside the survey's
scope, so the Hypothesis Testing model is not a theory of Jeffrey-type
sequence effects.}} \\
\multicolumn{3}{|p{\fullw}|}{\textbf{Quote.} \textit{Ortoleva 2024, Section 4.1, p. 559} \newline ``When new information $A$ arrives, individuals first test to
determine if they were using the right prior, using the threshold
$\epsilon$. If $A$ is not unlikely given $\pi$, that is, $\pi(A) > \epsilon$,
then $\pi$ is kept and updated following Bayes' rule \ldots when the new
information was unexpected given the prior, $\pi(A) \le \epsilon$, then the
prior is questioned''. \newline \textit{(part 4, Ortoleva) The departure is triggered by one unexpected
event, the AFTER's wording; the survey's restatement of Ortoleva (2012) is
the only source in hand.}} \\
\multicolumn{3}{|p{\fullw}|}{\textbf{Quote.} \textit{Cripps 2021, p. 9 (axioms named on pp. 7--9)} \newline ``Furthermore, symmetry implies that the order in which the
signals are revealed can be changed without an effect on the ultimate beliefs.
Hence, these axioms imply that reversing the order in which two signals arrive
has no effect on the ultimate beliefs.'' \newline \textit{(part 4, Cripps) His four axioms are Axiom 1 (Uninformativeness)
and Axiom 2 (Symmetry), p. 7, Axiom 3 (Divisibility), p. 8, and Axiom 4
(Non-Dogmatic), p. 9; order invariance follows from Symmetry and
Divisibility.}} \\
\multicolumn{3}{|p{\fullw}|}{\textbf{Theorem.} \textit{Cripps.lean, order\_invariance} \newline Let $U$ satisfy Symmetry and Divisibility, let $\mu$ have full
support, and let $x, y : \Theta \to (0,1)$. Writing $u(\mu,x)$ for the
first-signal update of $\mu$ on the binary experiment $(x, 1-x)$,
$u(u(\mu,x),y) = u(u(\mu,y),x)$. \newline \textit{(part 4, Cripps) The p. 9 remark as a theorem, with no Axiom 1
or 4 needed: two conditionally independent signals with fixed likelihoods
give the same posterior in either sequence.}} \\
\multicolumn{3}{|p{\fullw}|}{\textbf{Theorem.} \textit{Cripps.lean, composite\_AB\_eq\_bayes} \newline Let $\mu$ be a full-support belief on $A \times B$, $q > 0$ with
$\sum_a q_a = 1$ and $\sum_b r_b = 1$. Then
$J_B\big(J_A(\mu,q),r\big) = \mathrm{Bayes}(\mu, \ell)$ with
$\ell(a,b) = \dfrac{q_a}{\mu_A(a)} \cdot \dfrac{r_b}{\big(J_A(\mu,q)\big)_B(b)}$. \newline \textit{(part 4, Cripps) Under Proposition IMM's reading the composite
is one Bayes update, a divisible rule meeting all four axioms, whose second
likelihood is matched to the intermediate belief; the other sequence feeds a
different experiment, so no axiom fails.}} \\
\multicolumn{3}{|p{\fullw}|}{\textbf{Theorem.} \textit{Cripps.lean, rigid\_not\_uninformative} \newline For every $\varphi$, the rigid rule
$(\mu, p, s) \mapsto J_A\big(\mu, \varphi(p_s)\big)$ violates Axiom 1
(Uninformativeness): an experiment with $p_s(\theta) = p_s(\theta')$ for all
$\theta, \theta'$ still resets the $A$-marginal. \newline \textit{(part 4, Cripps) Under the other reading, delivered credence
fixed whatever the prior, the composite fails Uninformativeness (and
Non-Dogmatic), so ``of the four axioms only divisibility'' is wrong on either
reading.}} \\
\multicolumn{3}{|p{\fullw}|}{\textbf{Quote.} \textit{Pettigrew and Weisberg 2025, p. 3, Equation (1)} \newline ``But we'll see that commutativity instead favours upco, also
known as multiplicative pooling:
$P'(E) = \dfrac{P(E)\,Q(E)}{P(E)\,Q(E) + P(\overline{E})\,Q(\overline{E})}$. (1)'' \newline \textit{(part 5, Pettigrew-Weisberg) The pooling step is multiplicative
and takes the prior $P(E)$ as one of its two arguments.}} \\
\multicolumn{3}{|p{\fullw}|}{\textbf{Quote.} \textit{Pettigrew and Weisberg 2025, pp. 3--4} \newline ``Suppose we begin with $P$, fix some pooling rule $f$, and use the
following two-step procedure for responding to $Q$'s opinion about $E$.
\ldots Step 1. Apply pooling rule $f$ to $P(E)$ and $Q(E)$ to obtain
$P'(E)$'' (p. 3); ``Begin with the case where $P$ pools with $Q$ first. Step 1
of Jeffrey pooling combines $P(E) = 4/10$ with $Q(E) = 8/10$ via upco, to
yield $P'(E) = 8/11$'' (p. 4). \newline \textit{(part 5, Pettigrew-Weisberg) The prior is pooled with each
source's opinion; successive inputs are never pooled with each other.}} \\
\multicolumn{3}{|p{\fullw}|}{\textbf{Theorem.} \textit{PettigrewWeisberg.lean, upco\_eq\_self\_iff} \newline For $0 < p, q < 1$,
$\dfrac{pq}{pq + (1-p)(1-q)} = q \iff p = \tfrac12$. \newline \textit{(part 5, Pettigrew-Weisberg) Pooling the prior with a delivered
credence returns that credence only from an even prior, so their operation
is not the manuscript's marginal reset.}} \\
\multicolumn{3}{|p{\fullw}|}{\textbf{Theorem.} \textit{PettigrewWeisberg.lean, PB\_is\_upco\_pooling} \newline Let $P$ be regular on the cells $E_i \cap F_j$ and $x_i, y_j > 0$.
With the matched opinions $m_E(x)_i \propto x_i / P(E_i)$ and
$m_F(y)_j \propto y_j / P(F_j)$, Jeffrey pooling $P$ with upco on $E$ using
$m_E(x)$ and then on $F$ using $m_F(y)$, in either sequence, gives
$\PB(\omega) \propto P(\omega)\, \dfrac{x_{E(\omega)}}{P(E_{E(\omega)})}\,
\dfrac{y_{F(\omega)}}{P(F_{F(\omega)})}$. \newline \textit{(part 5, Pettigrew-Weisberg) Their commuting procedure, fed the
likelihoods matched to the prior, is Paper B's benchmark; sequence dependence
belongs to the delivered credences.}} \\ \hline
\multicolumn{3}{|l|}{} \\ \hline
\entryhead{3.3}{Related literature, Asch as evidence (was 3.B)\quad{\footnotesize\textit{(formerly E.6)}}}
 & While Bayes-factor updating predicts no sequence effect anywhere \citep{Wagner2002}, the belief-adjustment accounts of order effects \citep{HogarthEinhorn1992} focus on a single evaluative anchor. What the paper highlights is that such experimental analyses must also factor in the cross-attribute association -- thus requiring the elicitation of a full posterior table rather than a single rating produced by step-by-step or end-of-sequence response modes. & While Bayes-factor updating predicts no sequence effect anywhere \citep{Wagner2002}, the experimental tradition has measured the phenomenon entirely in marginals. \citet{HogarthEinhorn1992} track a single evaluative anchor. \citet{Asch1946} is more generous and more instructive. His Experiment~VI reports, for each of eighteen traits, the proportion of subjects choosing it over its opposite for the person described, separately for the two reading sequences. That the sequence effect differs in size across statistics is already visible there, although the differences Asch reports are all among marginals, which the paper treats alike. Between the sequences, \emph{restrained} moves from 64 to 9 per cent and \emph{good-looking} from 74 to 35, while \emph{serious} moves from 97 to 100, \emph{persistent} from 82 to 87 and \emph{reliable} from 84 to 91. Some statistics swing enormously and others scarcely move, and Asch attributes the difference to the content of the terms rather than to their position. Yet eighteen marginals separate full adoption from partial adoption no better than one. Separating them needs the rating of one trait before and after the second cue and against what that cue delivers (Proposition~\ref{prop:ADJ}), and the check-list asks whether a trait fits and never how strongly two traits are believed to go together beforehand, so the prior association cannot be formed from such data at all. What the paper adds is therefore not that more ratings are needed but that a second object must be elicited. Seeing the sequence effect requires only a marginal read in both sequences, which Asch has; locating the mechanism behind it requires the prior association as well, which he does not. \\ \hline
\multicolumn{3}{|p{\fullw}|}{\textbf{Why.} Plan 3.B with the audit applied. A11, the check-list is a forced choice between opposite pairs. A15, Asch does give an account of the uneven effects, by content. P13, what is never elicited is the prior association, not the joint. The ADJ clause now states what the recovery formula uses, three ratings of one marginal, not a marginal across values of $c$. ``amnestic updating'' becomes ``full adoption''. Colon removed. All eight percentages are correct and in the right columns.} \\
\multicolumn{3}{|p{\fullw}|}{\textbf{Evidence.} Asch.lean (quoted\_all, d\_neg\_exactly, stimulus\_disjoint\_checklist); Anchoring.lean (dampedB\_deviation, the identification).} \\ \hline
\multicolumn{3}{|p{\fullw}|}{\textbf{Grounds.}\par \textbf{Quote.} \textit{Asch 1946, p. 270 (Experiment VI)} \newline ``The following series are read, each to a different group: A. intelligent---industrious---impulsive---critical---stubborn---envious B. envious---stubborn---critical---impulsive---industrious---intelligent There were 34 subjects in Group A, 24 in Group B. The two series are identical with regard to their members, differing only in the order of succession of the latter.'' \newline \textit{Six stimulus terms read in either order; the eighteen traits of Table 7 are response items, none of which is a stimulus term (Asch.lean, stimulus\_disjoint\_checklist).}} \\
\multicolumn{3}{|p{\fullw}|}{\textbf{Quote.} \textit{Asch 1946, p. 262 (Check List I) and p. 271} \newline ``To this end we constructed a check list consisting of pairs of traits, mostly opposites. From each pair of terms in this list, which the reader will find reproduced in Table 1, the subject was instructed to select the one that was most in accordance with the view he had formed.'' (p. 262) ``The check-list data appearing in Table 7 furnish quantitative support for the conclusions drawn from the written sketches.'' (p. 271) \newline \textit{The check-list is a forced choice within a pair of opposites (audit A11); Table 7 reports, for each pair, the percentage choosing the listed member.}} \\
\multicolumn{3}{|p{\fullw}|}{\textbf{Computation.} \textit{Asch.lean, quoted\_all; literature/asch1946/sympy/check\_table7.py} \newline Table 7 (p. 271), Experiment VI, columns intelligent$\to$envious ($N = 34$) and envious$\to$intelligent ($N = 24$): restrained 64 / 9, good-looking 74 / 35, serious 97 / 100, persistent 82 / 87, reliable 84 / 91; also humorous 52 / 21, good-natured 18 / 0, important 85 / 90. Every quoted number is in the right column (proved by decide; script 16/16 PASS). \newline \textit{The four small moves are all of 7 points or less and all in the opposite direction to the large swings (Asch.lean, d\_neg\_exactly).}} \\
\multicolumn{3}{|p{\fullw}|}{\textbf{Quote.} \textit{Asch 1946, p. 264 (Experiment I)} \newline ``If we assume that the process of mutual influence took place in terms of the actual character of the qualities in question, it is not surprising that some will, by virtue of their content, remain unchanged.'' \newline \textit{Asch's account of uneven effects is by content; it is given for Experiment I, and he offers no trait-by-trait account for Experiment VI (audit A15).}} \\
\multicolumn{3}{|p{\fullw}|}{\textbf{Quote.} \textit{Asch 1946, pp. 271-272} \newline ``the first terms set up in most subjects a direction which then exerts a continuous effect on the latter terms.'' (p. 271) ``It is not the sheer temporal position of the item which is important as much as the functional relation of its content to the content of the items following it.'' (p. 272) \newline \textit{Content rather than position: the entry's ``rather than to their position''.}} \\
\multicolumn{3}{|p{\fullw}|}{\textbf{Theorem.} \textit{Anchoring.lean, dampedB\_deviation} \newline With $\mathrm{dampedB}(Q,r_0,\delta)$ the Jeffrey step on $B$ to $(1-\delta)\,Q(B{=}1) + \delta\,(1-r_0)$ and $Q(B{=}1) \neq 0$: $\mathrm{dampedB}(Q,r_0,\delta)(B{=}1) - (1-r_0) = (1-\delta)\,\big(Q(B{=}1) - (1-r_0)\big)$. Hence $\omega = \delta$ is fixed by three readings of one marginal, before the cue, after it, and the credence it delivers. \newline \textit{Marginals in two sequences, however many, do not fix omega; one marginal read before and after the second cue and against what it delivers does.}} \\
\multicolumn{3}{|p{\fullw}|}{\textbf{Quote.} \textit{Asch 1946, p. 262 (Technique) and p. 271 (design of the within-subject group)} \newline ``Following the reading, each subject wrote a brief sketch.'' (p. 262) ``A new group ($N{=}24$) heard Series B, wrote the free sketch, and immediately thereafter wrote the sketch in response to Series A.'' (p. 271) \newline \textit{Nothing is elicited before the series is read, and neither the sketch nor the forced-choice list asks how strongly two traits are believed to go together, so no prior association is collected (audit P13).}} \\ \hline
\multicolumn{3}{|l|}{} \\ \hline
\entryhead{3.5}{Related literature, Bohren-Imas-Rosenberg, after the identification paragraph (was 3.D)\quad{\footnotesize\textit{(formerly E.7)}}}
 & \textit{New paragraph(s) after the paragraph containing} ``contributes only a second-order loss in aggregate (Theorem~\ref{thm:LOS}).'' & A current formulation of the same problem is given by \citet{Bohren2019}, who distinguish discrimination arising from correct beliefs, from biased beliefs, and from preferences, and identify the source from how discrimination evolves along a history of evaluations. Their map from partiality to behaviour is Bayesian. A belief gap between groups is transmitted to evaluations, attenuated by the precision of the signal and along the history, and it vanishes as judgement becomes perfectly objective. Under the reading adopted here that map does not hold. A Jeffrey step sets the marginal of the attribute it addresses to the delivered credence, so an evaluator who takes the impression of quality last evaluates two workers alike whatever her prior beliefs about their groups, and the belief gap is silenced without any gain in objectivity. Under partial adoption with weight $\omega$ on the impression, a fraction $1-\omega$ of the belief gap survives (Proposition~\ref{prop:ADJ}). What credence-input updating does is therefore not to add a further source of discrimination to their three, but to relocate where partiality must sit in order to act. Lodged in the prior it is silenced, lodged in the impression it passes through untouched, and their framework has no parameter for the latter. \\ \hline
\multicolumn{3}{|p{\fullw}|}{\textbf{Why.} Plan 3.D with the audit applied. P17b, the belief gap is also attenuated along the history (their Proposition 2), so ``vanishes only as'' becomes ``vanishes as''. Colon removed. Goes directly after 3.4's paragraph, which carries the Heckman precedent that plan 3.C proposed (3.C is superseded by 3.4). Presupposes 5.2 and the Bohren2019 entry.} \\
\multicolumn{3}{|p{\fullw}|}{\textbf{Evidence.} check\_pinning\_kills\_partiality.py (7/7, D = 0 exactly for arbitrary group priors and covariances); Anchoring.lean dampedB\_deviation; BohrenImasRosenberg.lean (prop1\_decreasing, prop2\_decreasing, endo\_gap).} \\ \hline
\multicolumn{3}{|p{\fullw}|}{\textbf{Grounds.}\par \textbf{Quote.} \textit{Bohren, Imas and Rosenberg 2019, working paper p. 11 (Belief-Updating)} \newline ``The evaluator learns about the worker's ability from the evaluation history. Her posterior belief about ability is derived using Bayes rule, given her model of inference. She combines this updated belief about ability with the signal to learn about the quality of the current task, also using Bayes rule to form her posterior belief about quality.'' \newline \textit{Their map from partiality to behaviour is Bayesian by explicit assumption.}} \\
\multicolumn{3}{|p{\fullw}|}{\textbf{Quote.} \textit{Bohren, Imas and Rosenberg 2019, working paper pp. 11-12, eqs. (2) and (3)} \newline ``Then her optimal evaluation in period $t$ is $v_i(h_t,s_t,g) = \hat{E}_i[q_t \mid h_t,s_t,g] - c^i_g$. (2)'' (p. 11) ``Let $D_i(h,s) \equiv v_i(h,s,M) - v_i(h,s,F)$ (3) denote the difference between type $\theta_i$'s evaluation of a male and female worker conditional on observing history $h$ and signal $s$'' (p. 12). \newline \textit{With no taste ($c^i_g = 0$) discrimination is the gap in posterior expected quality at the same history and signal.}} \\
\multicolumn{3}{|p{\fullw}|}{\textbf{Quote.} \textit{Bohren, Imas and Rosenberg 2019, working paper p. 16, Proposition 1} \newline ``Proposition 1 (Subjectivity of Judgement). If the evaluator has belief-based partiality, initial discrimination is decreasing in the precision of the signal $\tau_\eta$ and otherwise, initial discrimination is constant with respect to $\tau_\eta$. As the signal becomes perfectly objective, $\tau_\eta \to \infty$, there is initial discrimination if and only if the evaluator has preference-based partiality.'' \newline \textit{The belief gap is attenuated by the precision of the signal and vanishes as judgement becomes perfectly objective.}} \\
\multicolumn{3}{|p{\fullw}|}{\textbf{Quote.} \textit{Bohren, Imas and Rosenberg 2019, working paper p. 18, Proposition 2} \newline ``Proposition 2 (Impossibility of Reversal). Suppose there is a single type of evaluator with belief-based partiality and no preference-based partiality. Then fixing an evaluation history, discrimination decreases across periods but never reverses.'' \newline \textit{The belief gap is also attenuated along the history, so the entry says ``vanishes as'' rather than ``vanishes only as'' (audit P17b).}} \\
\multicolumn{3}{|p{\fullw}|}{\textbf{Computation.} \textit{literature/bohren\_imas\_rosenberg2019/sympy/check\_pinning\_kills\_partiality.py, checks (3)-(5)} \newline Group $g$ has the $2{\times}2$ joint with $P_g(A{=}1) = a_g$, $P_g(B{=}1) = b_g$ and covariance $c_g$. A Jeffrey step on $A$ to $1-q_0$ followed by a Jeffrey step on $B$ to $1-r_0$ gives $P_M(B{=}1) - P_F(B{=}1) = 0$ identically in $(a_M,b_M,c_M,a_F,b_F,c_F,q_0,r_0)$, so $D = 0$ exactly with $c^i_g = 0$. With the $B$-cue read first and the $A$-cue last, $D \neq 0$. All checks pass. \newline \textit{An evaluator who takes the impression of quality last evaluates two workers alike whatever her group priors; the silencing is exact and route-dependent.}} \\
\multicolumn{3}{|p{\fullw}|}{\textbf{Theorem.} \textit{Anchoring.lean, dampedB\_deviation} \newline For $Q$ with $Q(B{=}1) \neq 0$: $\mathrm{dampedB}(Q,r_0,\delta)(B{=}1) - (1-r_0) = (1-\delta)\,\big(Q(B{=}1) - (1-r_0)\big)$. Applied to two groups meeting the same delivered credence $1-r_0$, the difference of their $B$-marginals after the cue is $(1-\delta)$ times the difference before it. \newline \textit{Under partial adoption with weight $\omega = \delta$ on the impression, a fraction $1-\omega$ of the belief gap survives.}} \\
\multicolumn{3}{|p{\fullw}|}{\textbf{Theorem.} \textit{BohrenImasRosenberg.lean, prop1\_decreasing, prop2\_decreasing} \newline With $\mathrm{eval}(\tau_q,\tau_\eta,\mu,c,s) = (\tau_q\mu + \tau_\eta s)/(\tau_q+\tau_\eta) - c$ and $D(\tau_\eta) = \mathrm{eval}(\tau_q,\tau_\eta,\hat\mu_M,0,s) - \mathrm{eval}(\tau_q,\tau_\eta,\hat\mu_F,c_F,s)$: if $\tau_q > 0$, $0 < \tau_{\eta,1} < \tau_{\eta,2}$ and $\hat\mu_F < \hat\mu_M$ then $D(\tau_{\eta,2}) < D(\tau_{\eta,1})$. Along any history $v : \mathbb{N} \to \mathbb{R}$, with $\tau_a,\tau_\varepsilon,\tau_\eta > 0$, $\hat\mu_F < \hat\mu_M$ and $c = 0$: $D_{n+1}(s) < D_n(s)$ for every $n$. \newline \textit{Their map transmits the belief gap attenuated by the signal precision and along the history; the paper's mechanism sits outside this map.}} \\ \hline
\multicolumn{3}{|l|}{} \\ \hline
\multicolumn{3}{|l|}{\rule{0pt}{3.2ex}\Large\textbf{4\quad Which statistics register the sequence}} \\ \hline\hline
\entryhead{4.2}{Section 4, the Foster-Greer-Thorbecke sentence\quad{\footnotesize\textit{(formerly C.13)}}}
 & This is the standard incidence-versus-intensity pairing of the measurement literature, where a headcount (a count of those past a threshold) and a mean shortfall (average of their distances from it) are first two members of a single parametrised family of indices \citep{FosterGreerThorbecke1984}. & The two decision statistics pair incidence with intensity, as do the first two poverty measures of \citet{FosterGreerThorbecke1984}, the headcount ratio, the share below a threshold, and the average shortfall below it normalised by the threshold, to which those above it contribute zero. \\ \hline
\multicolumn{3}{|p{\fullw}|}{\textbf{Why.} P\_0 is the headcount \emph{ratio}, a share, and P\_1 = (1/n) sum g\_i / z is normalised by the threshold and averaged over the whole population; ``average of their distances'' is the income-gap measure, which is not in the family (M15). FGT do not use ``incidence/intensity''; the AFTER does not attribute it to them.} \\
\multicolumn{3}{|p{\fullw}|}{\textbf{Evidence.} FGT.lean (\texttt{P0\_eq\_H}, \texttt{P1\_eq\_popMean\_normGap}, \texttt{I\_not\_decomposable}).} \\ \hline
\multicolumn{3}{|p{\fullw}|}{\textbf{Grounds.}\par \textbf{Quote.} \textit{Foster, Greer and Thorbecke 1984, p. 763, Equation (3)} \newline ``For each $\alpha \ge 0$, let $P_\alpha$ be defined by (3)
$P_\alpha(y; z) = \dfrac{1}{n} \sum_{i=1}^{q} \Big(\dfrac{g_i}{z}\Big)^{\alpha}$.'' \newline \textit{Every member divides by the population size $n$ and by the line
$z$; the sum runs over the $q$ poor only.}} \\
\multicolumn{3}{|p{\fullw}|}{\textbf{Quote.} \textit{Foster, Greer and Thorbecke 1984, p. 763} \newline ``The measure $P_0$ is simply the headcount ratio $H$, while
$P_1$ is $H \cdot I$, a renormalization of the income-gap measure. The
measure $P$ is obtained by setting $\alpha = 2$.'' \newline \textit{$P_0$ is a ratio, not a count, and $P_1$ is the income-gap
measure renormalised, not the ``average of their distances'' among the
poor; FGT never say ``incidence'' or ``intensity''.}} \\
\multicolumn{3}{|p{\fullw}|}{\textbf{Theorem.} \textit{FGT.lean, P1\_eq\_popMean\_normGap} \newline For any population $s$, incomes $y$ and line $z$,
$P_1(y; z) = \dfrac{1}{|s|} \sum_{i \in s} \dfrac{\max(z - y_i, 0)}{z}$. \newline \textit{$P_1$ is the whole-population mean of the normalised shortfall,
the non-poor contributing zero, which is the AFTER's description.}} \\
\multicolumn{3}{|p{\fullw}|}{\textbf{Theorem.} \textit{FGT.lean, I\_not\_decomposable} \newline With $z = 1$ and incomes $(0, \tfrac12, 2)$: $I = \tfrac34$ on the
whole population, $I = 1$ on $\{0\}$ and $I = \tfrac12$ on $\{1, 2\}$, and
$\tfrac34 \neq \tfrac13 \cdot 1 + \tfrac23 \cdot \tfrac12 = \tfrac23$. \newline \textit{The mean shortfall among the poor (FGT's $I$ times $z$) is not
additively decomposable with population-share weights, so it is not a member
of the decomposable family the manuscript cites.}} \\ \hline
\multicolumn{3}{|l|}{} \\ \hline
\multicolumn{3}{|l|}{\rule{0pt}{3.2ex}\Large\textbf{A\quad Appendix, Deferred proofs}} \\ \hline\hline
\entryhead{A.1}{Appendix, proof of Theorem LOS, the two Tao citations\quad{\footnotesize\textit{(formerly C.14)}}}
1/2 & Tonelli's theorem \citep[Corollary~1.7.23]{tao2011measure} gives & Tonelli's theorem \citep[Theorem~1.7.15]{tao2011measure}, applicable since the integrand is nonnegative, gives \\ \hline
2/2 & As $\mu$ is finite, continuity from above \citep[\S1.4]{tao2011measure} yields & As $\mu$ is finite and $B_c$ decreases as $c\downarrow0$, continuity from above \citep[Exercise~1.4.23(iii)]{tao2011measure}, applied along any sequence $c_n\downarrow0$, yields \\ \hline
\multicolumn{3}{|p{\fullw}|}{\textbf{Why.} Corollary 1.7.23 is Fubini-Tonelli; for the nonnegative integrand Tonelli (Theorem 1.7.15, or 1.7.18 for complete measures) is the exact reference. Continuity from above is Exercise 1.4.23(iii), stated for sequences; the proof takes c to 0 over a continuum (M16). Numbering is from the author's preprint of Tao; confirm against the printed book.} \\
\multicolumn{3}{|p{\fullw}|}{\textbf{Evidence.} Tao.lean (\texttt{thm\_1\_7\_15}, \texttt{ex\_1\_4\_23\_iii}, \texttt{measure\_band\_tendsto\_zero}, \texttt{los\_step4}).} \\ \hline
\multicolumn{3}{|p{\fullw}|}{\textbf{Grounds.}\par \textbf{Quote.} \textit{Tao 2011, Theorem 1.7.15, pre-p. 200} \newline ``Theorem 1.7.15 (Tonelli's theorem, incomplete version). Let
$(X, \mathcal{B}_X, \mu_X)$ and $(Y, \mathcal{B}_Y, \mu_Y)$ be $\sigma$-finite
measure spaces, and let $f : X \times Y \to [0, +\infty]$ be measurable with
respect to $\mathcal{B}_X \times \mathcal{B}_Y$. Then:'' \newline \textit{Tonelli is stated for nonnegative $f$ with no integrability
hypothesis, which is what Step 3 of the LOS proof uses.}} \\
\multicolumn{3}{|p{\fullw}|}{\textbf{Quote.} \textit{Tao 2011, Corollary 1.7.23, pre-p. 206} \newline ``Corollary 1.7.23 (Fubini-Tonelli theorem). Let
$(X, \mathcal{B}_X, \mu_X)$ and $(Y, \mathcal{B}_Y, \mu_Y)$ be complete
$\sigma$-finite measure spaces, and let $f : X \times Y \to \mathbf{C}$ be
measurable with respect to $\mathcal{B}_X \times \mathcal{B}_Y$. If
$\int_X \big(\int_Y |f(x,y)|\, d\mu_Y(y)\big)\, d\mu_X(x) < \infty$ \ldots
then $f$ is absolutely integrable''. \newline \textit{The corollary the manuscript cites is Fubini-Tonelli for
complex-valued $f$ under an absolute-integrability hypothesis, not the
nonnegative Tonelli statement.}} \\
\multicolumn{3}{|p{\fullw}|}{\textbf{Quote.} \textit{Tao 2011, Exercise 1.4.23(iii), pre-p. 93} \newline ``(iii) (Downwards monotone convergence) If $E_1 \supset E_2
\supset \ldots$ are $\mathcal{B}$-measurable, and $\mu(E_n) < \infty$ for at
least one $n$, then $\mu(\bigcap_{n=1}^{\infty} E_n) = \lim_{n \to \infty}
\mu(E_n) = \inf_n \mu(E_n)$.'' \newline \textit{Continuity from above is an exercise, stated for sequences and
with a finiteness hypothesis; the LOS proof lets $c \downarrow 0$ over a
continuum.}} \\
\multicolumn{3}{|p{\fullw}|}{\textbf{Theorem.} \textit{Tao.lean, measure\_band\_tendsto\_zero} \newline Let $\mu$ be a finite measure and $(B_c)_{c>0}$ measurable sets with
$B_{c'} \subseteq B_c$ whenever $0 < c' \le c$ and $\bigcap_{c > 0} B_c =
\emptyset$. Then $\mu(B_c) \to 0$ as $c \downarrow 0$. \newline \textit{The real-parameter passage Step 4 needs, proved from the
sequential statement since the neighbourhood filter of $0^{+}$ is countably
generated; the AFTER applies 1.4.23(iii) along any $c_n \downarrow 0$.}} \\ \hline
\multicolumn{3}{|l|}{} \\ \hline
\multicolumn{3}{|l|}{\rule{0pt}{3.2ex}\Large\textbf{B\quad Back matter and bibliography}} \\ \hline\hline
\entryhead{B.1}{Back matter, AI declaration (was B.B)\quad{\footnotesize\textit{(formerly E.13)}}}
 & Proposition~\ref{prop:PRO} (uniqueness of the protected statistic), which was then independently re-derived & Propositions~\ref{prop:PRO} (uniqueness of the protected statistic), \ref{prop:ORD} (between-sequence contrast), \ref{prop:ADJ} (adoption weight), \ref{prop:LAD} (partial adoption) and~\ref{prop:FAC} (factor inputs), which were then independently re-derived \\ \hline
\multicolumn{3}{|p{\fullw}|}{\textbf{Why.} Plan B.B. Propositions ORD, ADJ, LAD and FAC have the same provenance as PRO, so the declaration names them once 5.2 and 6.3 are applied.} \\
\multicolumn{3}{|p{\fullw}|}{\textbf{Evidence.} None needed.} \\ \hline
\multicolumn{3}{|p{\fullw}|}{\textbf{Grounds.}\par \textbf{Quote.} \textit{PAPER\_B\_MANUSCRIPT.tex, lines 1220-1228 (author's draft), Declaration of Generative AI} \newline ``During the preparation of this work the author used Anthropic's Claude (large
language model) in order to draft and refine manuscript prose; propose, prove, and
cross-verify Proposition~\ref{prop:PRO} (uniqueness of the protected statistic), which
was then independently re-derived and machine-verified by the author; propose
bibliographic sources for positioning against neighbouring literatures; and structure
the manuscript to conform to the target journal's house style.'' \newline \textit{Propositions ORD and ADJ were proposed, proved and cross-verified the same way (PropORD.lean, Anchoring.lean, verify\_ORD.py, check\_zero\_slope\_identification.py), so the declaration names them with PRO once 5.2 is applied.}} \\ \hline
\multicolumn{3}{|l|}{} \\ \hline
\entryhead{B.2}{Bibliography entries the Section E texts need (was B.A)\quad{\footnotesize\textit{(formerly E.14)}}}
 & \multicolumn{2}{p{24.6cm}|}{\textit{No manuscript text; see Why.}} \\ \hline
\multicolumn{3}{|p{\fullw}|}{\textbf{Why.} Needed by 1.5, 3.1, 3.5 (Bohren2019, Doring1999, Garber1980) and 3.4 (Heckman1998, Bohren2019). Hawthorne2004 is already in bibliography.bib. All six proposed entries are in notes/manuscript\_corrections\_extra.bib and move to bibliography.bib on approval. Zhao-Osherson is cited only in the introduction (author's draft) and is already in the bibliography.} \\
\multicolumn{3}{|p{\fullw}|}{\textbf{Evidence.} Bibliographic details verified against the Drive copies where printed (Doring: Phil. Sci. 66 (Proceedings) S379-S389; Garber: 47(1) 142-145; Heckman: JEP 12(2) 101-116); Bohren et al.'s AER pages are from the reference lists of later papers, the Drive copy being the working paper.} \\ \hline
\multicolumn{3}{|l|}{} \\ \hline
\end{longtable}
}

\section*{Bibliography entries the AFTER texts need (B.2)}
\texttt{Hawthorne2004} and \texttt{ZhaoOsherson2010} are already in \texttt{bibliography.bib}.
Six more are needed by 1.6, 2.4, 3.4, 1.5, 3.1 and 3.5 (Jeffrey 2004, Benjamin et al.\ 2019, Heckman 1998,
Bohren et al.\ 2019, D\"oring 1999, Garber 1980); they are in
\texttt{notes/manuscript\_corrections\_extra.bib} so that this document renders them, and move to
\texttt{bibliography.bib} on approval. The Drive copy of Jeffrey (2004) is the November 2002 draft,
so 1.6 and 2.4 cite the chapter only.

\section*{Where each entry of the archived plan (notes/manuscript\_change\_plan\_asof\_2026-09-30.md) now lives}
{\small
\begin{longtable}{|p{2.6cm}|p{5.2cm}|p{17.4cm}|}
\hline
\textbf{Entry} & \textbf{Subject} & \textbf{Status and what it needs} \\ \hline\endhead
0.A & Abstract & Superseded by the author's rewrite and 0.1. \\ \hline
1.A, 1.B, 1.D, 1.E & Introduction & Superseded by the author's rewrite; what survives is corrected in 1.1, 1.2, 1.3, 1.4, 1.6, 1.7 and 1.8. The ORD sentence of 1.E is still unapplied and presupposes 5.2. \\ \hline
1.A2 & Two-channel prelude & Applied by the author; corrected by 1.3. \\ \hline
1.C & Why sequence matters beyond one judgement & Re-issued as 1.5. \\ \hline
2.A, 2.B, 2.C & Impression; worked example; soft cues & Re-issued as 2.2, 2.5 (extended with the soft-versus-hard passage), 2.6. \\ \hline
3.A, 3.B & Commutativity literature; Asch & Re-issued as 3.1, 3.3. \\ \hline
3.C & Heckman & Superseded by the identification paragraph appended in 3.4, with P17a applied. \\ \hline
3.D & Bohren-Imas-Rosenberg & Re-issued as 3.5. \\ \hline
4.A, 5.A & Section titles and openings & Re-issued as 4.1, 5.1. \\ \hline
5.B & Propositions ORD and ADJ & ORD re-issued as 5.2; ADJ moved, with LAD and FAC, to 6.3. \\ \hline
6.A, 6.B & Rival mechanisms; two channels & Re-issued as 6.2 and 6.3 (6.B's text pulled in from interior\textbackslash{}\_omega.tex). \\ \hline
B.A, B.B & Bibliography; AI declaration & Re-issued as B.2, B.1. \\ \hline
W.A & Input/output tables & Withdrawn by the author 2026-09-21; kept in notes/empirical\textbackslash{}\_analytics.tex. \\ \hline
\end{longtable}
}

\section*{Old and new entry numbers}
{\small
\begin{longtable}{|p{2.2cm}|p{2.2cm}|p{9cm}|p{4cm}|}
\hline
\textbf{Former} & \textbf{Now} & \textbf{Section} & \textbf{Status} \\ \hline\endhead
C.1 & 0.1 & Abstract & pending \\ \hline
C.2 & 1.1 & Introduction & pending \\ \hline
C.3 & 1.2 & Introduction & pending \\ \hline
C.4 & 1.3 & Introduction & pending \\ \hline
C.5 & 1.4 & Introduction & pending \\ \hline
C.6 & 1.6 & Introduction & pending \\ \hline
C.7 & 1.7 & Introduction & pending \\ \hline
C.8 & 1.8 & Introduction & pending \\ \hline
C.9 & 2.1 & Setup & applied at db3a2f41 \\ \hline
C.9w & 2.3 & Setup & applied at db3a2f41 \\ \hline
C.10 & 2.4 & Setup & applied at db3a2f41 \\ \hline
C.11 & 3.2 & Related literature & pending \\ \hline
C.11b & 3.4 & Related literature & applied at db3a2f41 \\ \hline
C.13 & 4.2 & Which statistics register the sequence & pending \\ \hline
C.14 & A.1 & Appendix, Deferred proofs & pending \\ \hline
E.1 & 1.5 & Introduction & pending \\ \hline
E.2 & 2.2 & Setup & applied at db3a2f41 \\ \hline
E.3 & 2.5 & Setup & applied at db3a2f41 \\ \hline
E.4 & 2.6 & Setup & pending \\ \hline
E.5 & 3.1 & Related literature & applied at db3a2f41 \\ \hline
E.6 & 3.3 & Related literature & pending \\ \hline
E.7 & 3.5 & Related literature & pending \\ \hline
E.8 & 4.1 & Which statistics register the sequence & applied at db3a2f41 \\ \hline
E.9 & 5.1 & Averaging over sequences does not change the classification & applied at db3a2f41 \\ \hline
E.10 & 5.2 & Averaging over sequences does not change the classification & applied at db3a2f41 \\ \hline
E.11 & 6.2 & Scope and robustness & applied at db3a2f41 \\ \hline
E.13 & B.1 & Back matter and bibliography & pending \\ \hline
E.14 & B.2 & Back matter and bibliography & pending \\ \hline
E.15 & 6.3 & Scope and robustness & applied at db3a2f41 \\ \hline
E.15h & 6.1 & Scope and robustness & applied at db3a2f41 \\ \hline
E.15r & 1.9 & Introduction & pending \\ \hline
C.12 & B.2 & folded in & \\ \hline
C.15 & 3.4 & folded in & \\ \hline
E.12 & 6.3 & folded in & \\ \hline
\end{longtable}
}

\bibliographystyle{plainnat}
\bibliography{../bibliography,manuscript_corrections_extra}
\end{document}
